\documentclass[11pt,letterpaper]{article}



\usepackage[a4paper,margin=23mm,includeheadfoot,headheight=15pt,headsep=15pt,footskip=22pt]{geometry}
\usepackage{setspace,needspace,titlesec,etoolbox}
\usepackage[nopatch=footnote]{microtype}
\usepackage{amsmath,amssymb,amsthm,mathtools}
\usepackage{aliascnt}
\usepackage{booktabs,longtable,array,tabularx}
\usepackage{xcolor}
\usepackage{enumitem}
\usepackage{fancyhdr}
\usepackage{xurl}
\usepackage{hyperref}
\usepackage[nameinlink,noabbrev]{cleveref}
\usepackage{bookmark,graphicx,caption,tikz,multicol}
\usetikzlibrary{arrows.meta,decorations.pathreplacing,patterns,calc}

\allowdisplaybreaks
\hyphenpenalty=1000
\exhyphenpenalty=1000
\tolerance=2000
\emergencystretch=3em

% Keep the release PDF deterministic across clean rebuilds.
\ifdefined\pdfinfoomitdate
  \pdfinfoomitdate=1
\fi
\ifdefined\pdftrailerid
  \pdftrailerid{}
\fi
\ifdefined\pdfsuppressptexinfo
  \pdfsuppressptexinfo=15
\fi

\setlist[enumerate]{topsep=4pt,itemsep=2pt}
\setlist[itemize]{topsep=4pt,itemsep=2pt}

\newtheorem{theorem}{Theorem}[section]
\newaliascnt{proposition}{theorem}
\newtheorem{proposition}[proposition]{Proposition}
\aliascntresetthe{proposition}
\newaliascnt{lemma}{theorem}
\newtheorem{lemma}[lemma]{Lemma}
\aliascntresetthe{lemma}
\newaliascnt{corollary}{theorem}
\newtheorem{corollary}[corollary]{Corollary}
\aliascntresetthe{corollary}
\theoremstyle{definition}
\newaliascnt{definition}{theorem}
\newtheorem{definition}[definition]{Definition}
\aliascntresetthe{definition}
\newaliascnt{example}{theorem}
\newtheorem{example}[example]{Example}
\aliascntresetthe{example}
\newaliascnt{standardfact}{theorem}
\newtheorem{standardfact}[standardfact]{Standard Fact}
\aliascntresetthe{standardfact}
\theoremstyle{remark}
\newaliascnt{remark}{theorem}
\newtheorem{remark}[remark]{Remark}
\aliascntresetthe{remark}
\newaliascnt{boundary}{theorem}
\newtheorem{boundary}[boundary]{Boundary}
\aliascntresetthe{boundary}

\crefname{theorem}{theorem}{theorems}
\crefname{proposition}{proposition}{propositions}
\crefname{lemma}{lemma}{lemmas}
\crefname{corollary}{corollary}{corollaries}
\crefname{definition}{definition}{definitions}
\crefname{example}{example}{examples}
\crefname{standardfact}{standard fact}{standard facts}
\crefname{remark}{remark}{remarks}
\crefname{boundary}{boundary}{boundaries}

\newcommand{\CC}{\mathbb C}
\newcommand{\NN}{\mathbb N}
\newcommand{\ZZ}{\mathbb Z}
\newcommand{\Bcal}{\mathcal B}
\newcommand{\Ecal}{\mathcal E}
\newcommand{\Pcal}{\mathcal P}
\newcommand{\Rcal}{\mathcal R}
\newcommand{\Zcal}{\mathcal Z}
\newcommand{\Dfun}{\mathsf D}
\newcommand{\Efun}{\mathsf E}
\newcommand{\Sfun}{\mathsf S}
\newcommand{\Tfun}{\mathsf T}
\newcommand{\Cfun}{\mathsf C}
\newcommand{\Thetafun}{\symup{\Theta}}
\newcommand{\Spec}{\operatorname{Spec}}
\newcommand{\ord}{\operatorname{ord}}
\newcommand{\type}{\operatorname{type}}
\newcommand{\Log}{\operatorname{Log}}
\newcommand{\Res}{\operatorname{Res}}
\newcommand{\dist}{\operatorname{dist}}
\newcommand{\sech}{\operatorname{sech}}
\newcommand{\supp}{\operatorname{supp}}
\newcommand{\id}{\mathrm{id}}
\newcommand{\tr}{\operatorname{tr}}
\newcommand{\Fdet}{\operatorname{det}_{\!F}}
\newcommand{\Span}{\operatorname{span}}
\newcommand{\leanname}[1]{\nolinkurl{#1}}
\newcolumntype{Y}{>{\raggedright\arraybackslash}X}

\definecolor{leanblue}{RGB}{20,72,135}
\definecolor{accent}{HTML}{137668}
\definecolor{muted}{HTML}{536568}
\frenchspacing\clubpenalty=10000\widowpenalty=10000\raggedbottom
\titleformat{\section}{\sffamily\bfseries\large\color{accent}}{\thesection}{.7em}{}
\titleformat{\subsection}{\sffamily\bfseries\normalsize}{\thesubsection}{.7em}{}
\titlespacing*{\section}{0pt}{1.4\baselineskip}{.55\baselineskip}
\pretocmd{\section}{\Needspace{7\baselineskip}}{}{}
\pretocmd{\subsection}{\Needspace{5\baselineskip}}{}{}
\setstretch{1.055}\setlength{\parskip}{.3em}\setlength{\parindent}{1.35em}
\definecolor{statusgray}{RGB}{86,86,86}

\newcommand{\statusleanfour}[1]{%
  \par\smallskip\noindent
  {\color{leanblue}\textbf{Released formal contract: Paper IV v1.0.0.}}\par
  \begingroup\small\raggedright\leanname{#1}\par\endgroup\smallskip}
\newcommand{\statusleanprior}[1]{%
  \par\smallskip\noindent
  {\color{leanblue}\textbf{Inherited formal contract: earlier papers.}}\par
  \begingroup\small\raggedright\leanname{#1}\par\endgroup\smallskip}

\title{\textbf{Projector Resolved Seed Calculus IV}\\[0.45em]
\large Cyclic Analytic Functions, Reciprocal Families, and Spectral Zero Data}
\author{Komorebi P.}
\date{Modern Monograph, 8 October 2026}

\hypersetup{
  pdftitle={Projector-Resolved Seed Calculus IV: Cyclic Analytic Functions, Reciprocal Families, and Spectral Zero Data},
  pdfauthor={Komorebi P.},
  pdfsubject={Cyclic exponential sections, reciprocal families, zero geometry, and spectral zero data},
  colorlinks=true,
  linkcolor=accent,
  citecolor=accent,
  urlcolor=accent
}

\pagestyle{fancy}
\fancyhf{}
\fancyhead[L]{\sffamily\fontsize{8}{10}\selectfont\color{muted}KOMOREBI MATHEMATICS / PAPER IV}
\fancyhead[R]{\sffamily\fontsize{8}{10}\selectfont\color{muted}TRADITIONAL PAPER}
\fancyfoot[L]{\sffamily\fontsize{8}{10}\selectfont\color{muted}Projector-resolved seed calculus}
\fancyfoot[R]{\sffamily\small\color{accent}\thepage}


% Operation Komorebi: shared publication profile.
\usepackage{komorebi-monograph}
\renewcommand{\KomorebiTitle}{Projector Resolved Seed Calculus IV}
\renewcommand{\KomorebiSubtitle}{Cyclic analytic functions, reciprocal families,\\and spectral zero data}
\renewcommand{\KomorebiSeries}{Paper IV}
\renewcommand{\KomorebiKind}{Modern Monograph}
\renewcommand{\KomorebiEdition}{8 October 2026}
\renewcommand{\KomorebiShortTitle}{Paper IV}
\renewcommand{\KomorebiTitleSize}{31}
\renewcommand{\KomorebiTitleLeading}{36}
\renewcommand{\KomorebiProfile}{monograph}
\renewcommand{\KomorebiOrientation}{portrait}
\renewcommand{\KomorebiCoverTitle}{Projector Resolved\\Seed Calculus IV}
\renewcommand{\KomorebiCoverBlurb}{One cyclic family connects exact coefficient selection, differential transport, reciprocal polynomials, and the geometry of zeros.}
\KomorebiSetup
\graphicspath{{figures/}}
\providecommand{\OO}{\mathcal O}
\providecommand{\FF}{\mathbb F}\providecommand{\QQ}{\mathbb Q}\providecommand{\ZZ}{\mathbb Z}\providecommand{\RR}{\mathbb R}
\providecommand{\Spec}{\operatorname{Spec}}\providecommand{\Supp}{\operatorname{Supp}}\providecommand{\Fitt}{\operatorname{Fitt}}
\providecommand{\Ann}{\operatorname{Ann}}\providecommand{\coker}{\operatorname{coker}}\providecommand{\im}{\operatorname{im}}
\providecommand{\rank}{\operatorname{rank}}\providecommand{\Tr}{\operatorname{Tr}}\providecommand{\id}{\operatorname{id}}
\providecommand{\whot}{\mathbin{\widehat\otimes}}\providecommand{\m}{\mathfrak m}\providecommand{\p}{\mathfrak p}
\providecommand{\head}[1]{\par\vspace{2pt}{\sffamily\bfseries\fontsize{10}{12}\selectfont #1}\par\vspace{1pt}}
\providecommand{\platecaption}[1]{\par\noindent{\footnotesize\color{black!70}#1\par}}
\newtcolorbox{insight}{colback=accent!3,colframe=accent,boxrule=0pt,leftrule=1.1pt,arc=0pt,left=10pt,right=10pt,top=5pt,bottom=5pt,fontupper=\fontsize{10}{12}\selectfont,before skip=6pt,after skip=5pt}
\providecommand{\takeaway}[1]{\begin{insight}\textbf{Read this as:}\ #1\end{insight}}
\providecommand{\vfig}[2][.92]{\par\begingroup\centering\resizebox{#1\linewidth}{!}{\input{visual_graphics/#2.tex}}\par\endgroup\vspace{2pt}}

\newif\ifstarted\startedfalse
\providecommand{\vpage}[4]{\ifstarted\clearpage\else\startedtrue\fi
\pdfbookmark[1]{V#1. #2}{plate#1}\hypertarget{vpage#1}{}\label{visual:#1}
{\sffamily\fontsize{9}{11}\selectfont\bfseries\color{accent}V#1\quad /\quad #3\par}\vspace{5pt}
{\sffamily\bfseries\fontsize{19}{22}\selectfont #2\par}\vspace{5pt}
{\sffamily\fontsize{8.7}{11}\selectfont\color{muted} #4\par}\vspace{6pt}}
\providecommand{\pagerefV}[1]{\hyperlink{vpage#1}{V\,--\,#1}}

\providecommand{\CC}{\mathbb C}
\providecommand{\kk}{\mathbb k}
\renewcommand{\AA}{\mathcal A}
\providecommand{\BB}{\mathcal B}
\providecommand{\Hull}{\operatorname{Hull}}
\providecommand{\srank}{\operatorname{srank}}
\providecommand{\Span}{\operatorname{span}}
\providecommand{\End}{\operatorname{End}}
\providecommand{\diag}{\operatorname{diag}}
\providecommand{\Ker}{\operatorname{ker}}
\providecommand{\norm}[1]{\left\lVert#1\right\rVert}
\providecommand{\mat}[1]{\begin{pmatrix}#1\end{pmatrix}}

\providecommand{\Ob}{\operatorname{Ob}}
\providecommand{\supp}{\operatorname{supp}}
\providecommand{\src}{\operatorname{src}}\providecommand{\tgt}{\operatorname{tgt}}


\providecommand{\captiontext}[1]{\platecaption{#1}}

\captionsetup{font={small,sf},labelfont={bf,sf},labelsep=period}
\newcommand{\modelfigure}[3]{\begin{figure}[htbp]\centering\includegraphics[width=.98\linewidth]{#1.pdf}\caption{#2}\label{#3}\end{figure}}
\usepackage{titlesec}
\input{core-layout.tex}

% Publication-design overrides. Mathematical bodies retain their source text.
\CoreLayout
\colorlet{accent}{pubTeal}\colorlet{teal}{pubTeal}\colorlet{ink}{pubInk}\colorlet{muted}{pubMuted}\colorlet{pale}{pubWash}\colorlet{light}{pubWash}
\makeatletter
\@ifundefined{KomorebiCover}{}{\renewcommand{\KomorebiCover}{%
\clearpage\hypersetup{pageanchor=false}\thispagestyle{empty}\begingroup\setlength{\parindent}{0pt}%
{\sffamily\fontsize{8.5}{11}\selectfont\color{pubTeal}REALIZATION CALCULUS\par}%
\vspace{27pt}%
{\sffamily\fontsize{9}{12}\selectfont\color{pubTeal}\MakeUppercase{\KomorebiSeries}\par}%
\vspace{22pt}%
{\normalfont\fontsize{29}{35}\selectfont\color{pubNavy}\raggedright\KomorebiTitle\par}%
\vspace{17pt}%
{\normalfont\fontsize{17}{23}\selectfont\color{pubMuted}\raggedright\KomorebiSubtitle\par}%
\vspace{22pt}\noindent\textcolor{pubCopper}{\rule{42pt}{.8pt}}\par%
\vspace{20pt}%
\@ifundefined{KomorebiCoverBlurb}{}{\normalfont\fontsize{11.4}{16.2}\selectfont\color{pubMuted}\raggedright\KomorebiCoverBlurb\par}%
\vfill%
{\normalfont\fontsize{14}{18}\selectfont\color{pubInk}Komorebi P.\par}%
\vspace{9pt}{\sffamily\fontsize{9}{13}\selectfont\color{pubMuted}\KomorebiKind\par\KomorebiEdition\par}%
\endgroup\clearpage\pagenumbering{roman}\hypersetup{pageanchor=true}\color{pubInk}\normalfont\pagestyle{empty}}}
\makeatother

\begin{document}
\KomorebiCover
\KomorebiOpeningNote{Begin with cyclic projection and the signed differential cycle. The reciprocal chapters then build parameter polynomials, Appell systems, and theta polynomials from a unit series. The second half follows the same denominator through large zero geometry, canonical products, and zero zeta functions.

The visual atlas at the end of this volume supplies twenty-six worked plates. The separate Lean guide follows the released declarations and their hypotheses. All sixty-eight proof bodies and ninety-five numbered records of the selected manuscript are retained.}

\begin{abstract}
Fix a positive integer $N$. Root of unity projection of the exponential
produces residue-class sections that, for $\lambda\ne0$, form a basis for
$y^{(N)}=\lambda^N y$. The signed sections for $y^{(N)}=-y$ yield the
denominator $\Dfun_N$ and its reduced form $\Delta_N$. We derive reciprocal
coefficient polynomials, Bell and Newton recurrences, an additive parameter
law, and residue section Appell and theta families. For $N\geq2$,
$\Delta_N$ has order $1/N$. We specialize classical two term zero theory for
exponential sums to the regular exponent polygon. In this equal-weight family,
conjugation symmetry combined with unique near-model localization places all
sufficiently large zeros exactly on the symmetry rays, and direct expansion
gives their leading displacement and derivative sign. The classical perimeter
law yields zero-counting asymptotics with explicit leading constants and
bounded error. These results give genus zero products, the standard diagonal
Fredholm representation, zero-power-sum formulas for reciprocal and
logarithmic coefficients, nearest-zero coefficient asymptotics, and
meromorphic continuation of the symmetry-orbit zero zeta by comparison with
the Hurwitz zeta function. Complex powers are interpreted either formally or
through a fixed logarithm on a zero free domain. The paper treats only this
cyclic analytic family. Arithmetic interpolation, general nonlinear trees,
general resolvent theory, and abstract completion are reserved for Papers IX,
III, VI, and X.
\end{abstract}

\begin{readerroadmap}[title={Formal scope and evidence}]
The preserved Paper IV v1.0.0 release records full scoped Lean coverage of the sixty-eight theorem like manuscript results under their original hypotheses. This volume changes presentation and adds worked explanations. It does not create a new formal release.

Imported standard fact 11.1 is used in the ordinary manuscript growth argument. The Lean growth-order proofs establish that step separately. Open questions, branch choices, multiplicities, and the hypotheses of coefficient asymptotics remain explicit. Current finite computations support the worked models. The compiler and kernel evidence remains the preserved September release evidence.
\end{readerroadmap}

\clearpage
\tableofcontents
\listoffigures
\clearpage\pagenumbering{arabic}

\section{Introduction}\label{sec:introduction}

For a power series
\[
F(z)=\sum_{m\geq 0}a_mz^m
\]
and a positive integer $N$, the coefficients may be separated according to the residue of $m$ modulo $N$.  Root of unity averaging gives exact projection operators for this separation.  This elementary construction is usually called a multisection or a root filter.  It appears throughout the classical theory of generating functions and special functions.  Background on multisections, hypergeometric functions, and related coefficient systems may be found in \cite{Rainville1960,AndrewsAskeyRoy1999,Comtet1974}.

Papers I and II developed typed seed grammars, exact projector packets, and
finite residue-projector algebra, while Paper III treated typed edge operators
and exact coordinate transport \cite{PaperI2026,PaperII2026,PaperIII2026}.
The root filter used below is the finite cyclic scalar specialization of that
earlier projector calculus. The signed companion matrix is likewise a concrete
cyclic instance, not a replacement for Paper III's general typed-edge theory.

This volume follows the consequences of applying the cyclic projections to the exponential function.  The projected exponentials form a finite differential system.  After one sign is inserted around the derivative cycle, they form a basis for
\[
y^{(N)}=-y.
\]
The section with initial value one is
\[
\Dfun_N(z)=\sum_{n\geq 0}(-1)^n\frac{z^{Nn}}{(Nn)!}.
\]
Its reciprocal powers produce polynomial coefficient families. Its logarithmic
derivative produces a scalar higher Riccati identity. The same reciprocal
construction applied to every exponential section gives normalized
residue section functions and ordinary Appell polynomial sequences. These
cyclic exponential sections and their matrix addition laws belong to the
generalized-hyperbolic-function literature
\cite{Schwaiger1992,Muldoon2005}. Appell and binomial-type sequences are
discussed in \cite{Roman1984,Comtet1974}.

The second half of the paper concerns the zeros of $\Dfun_N$.  The finite exponential sum
\[
\Dfun_N(z)=\frac1N\sum_{\eta^N=-1}e^{\eta z}
\]
has a regular polygon of exponents.  The general relation between an exponent polygon and the zero strips of an exponential sum goes back to P\'olya and Langer.  Later developments include the work of Moreno and the treatment in Levin's monograph \cite{Polya1920,Langer1931,Moreno1973,Levin1980}. Bleher and Mallison give a particularly close general result: for convex-polygon exponential sums they locate all sufficiently large zeros in two term strips and prove exponential localization and simplicity \cite{BleherMallison2006}. The regular polygon and conjugation symmetry of the present family specialize that theory to $N$ exact rays. We give a self-contained specialization because its notation is used throughout the product and zeta theory, and we compute the leading displacement term for this family. The resulting counting formulas have explicit leading constants and bounded error.

The rotational symmetry is more useful after the change of variable $w=z^N$.  There is an entire function
\[
\Delta_N(w)=\sum_{n\geq 0}(-1)^n\frac{w^n}{(Nn)!}
\]
with
\[
\Dfun_N(z)=\Delta_N(z^N).
\]
For $N\geq2$, the order of $\Delta_N$ is $1/N$.  Its Hadamard product therefore has genus zero.  The logarithm of this product identifies the transfer tangent coefficients with power sums of the reduced zeros.  The large zero formula also compares the reduced zero zeta with a Hurwitz zeta function.  Standard material on entire functions, canonical products, and exponential sums appears in \cite{Boas1954,Levin1980,Titchmarsh1939}.  We use singularity analysis in the form described in \cite{FlajoletSedgewick2009} and the classical continuation of the Hurwitz zeta function from \cite{Apostol1976}.

The cyclic exponential sections and their matrix addition laws have a
classical literature. The tools used below include Appell and binomial type
sequences, canonical products, singularity analysis, and Fredholm
determinants. These tools are also standard. Our aim is to organize this symmetric
specialization and link its reciprocal coefficients to reduced-zero data. No
historical priority is asserted for the component constructions.



\section{A map of the construction}\label{sec:map}

\begin{keyidea}
The same object is read in different coordinates. Selecting residues makes its Taylor coefficients visible. A signed shift makes the differential system finite. Passing to $w=z^N$ removes rotational repetition from its zeros. The formulas below identify these descriptions under their stated hypotheses.
\end{keyidea}

Two variables appear throughout the paper.  The variable $z$ carries the differential equation and the rotational symmetry.  The reduced variable
\[
w=z^N
\]
records one representative from each rotational orbit.  Confusing these variables hides the main simplification.  The denominator $\Dfun_N$ has order one as a function of $z$, while $\Delta_N$ has order $1/N$ as a function of $w$.

The analytic objects are introduced in the following order.  First,
\[
\Efun_{N,j}(z)
=
\sum_{n\geq0}\frac{z^{Nn+j}}{(Nn+j)!}
\]
is the residue $j$ part of the exponential.  Second,
\[
\Dfun_{N,j}(z)
=
\sum_{n\geq0}(-1)^n\frac{z^{Nn+j}}{(Nn+j)!}
\]
is the signed basis for $y^{(N)}=-y$.  The section $\Dfun_N=\Dfun_{N,0}$ is the denominator used in the reciprocal family
\[
\Sfun_N^{(r)}=\Dfun_N^{-r}.
\]
Set $\Tfun_N=-\Dfun_N^{\prime}/\Dfun_N$. Equivalently, $\Tfun_N$ is the logarithmic derivative of $\Sfun_N^{(1)}$.

A different reciprocal family comes from normalizing each exponential section at the origin.  The function
\[
\Cfun_{N,j}(z)
=
\frac{z^j/j!}{\Efun_{N,j}(z)}
\]
is a unit series in $z^N$.  Its powers, after multiplication by $e^{xz}$, give Appell polynomials.  The theta family couples a scaled argument with a reciprocal order.

The zero theory begins only after these algebraic relations have been established.  The factorization
\[
\Dfun_N(z)=\Delta_N(z^N)
\]
turns every orbit of $N$ zeros in the $z$ plane into one zero in the $w$ plane.  The large zeros in the $z$ plane are located by two dominant exponentials.  The zeros in the $w$ plane then give a genus zero product.  The logarithm of that product returns to the coefficient theory through
\[
\text{zero power sums}
=
\text{logarithmic coefficients}.
\]
This equality is the bridge between the first and second halves of the paper.

All powers with a complex parameter are first defined as formal series at the origin.  To represent this germ analytically, use a simply connected zero free region containing the origin and the logarithm normalized to vanish there.  On other regions, specify a branch separately.  Integer reciprocal powers are globally meromorphic and do not require a branch choice.

The main statements fall into five groups.  The finite algebra model in \cref{sec:finite_algebra} explains the derivative cycles and addition laws through one exponential identity.  The coefficient theory in \cref{sec:reciprocal_powers} gives two complementary recurrences, one from the original denominator and one from its logarithm.  The Appell theory in \cref{sec:residue_section} supplies both formal differential-operator formulas and a convolution law in the space and order parameters.  The zero analysis in \cref{sec:zero_geometry} specializes classical two term exponential-sum theory, uses symmetry to place the uniquely localized large zeros on the rays, and computes a leading correction.  The product theory in \cref{sec:hadamard} gives the standard diagonal Fredholm encoding of the genus zero product.

The volume keeps formal and analytic statements separate.  A formal identity needs only the unit constant term of the relevant series.  A holomorphic complex power also needs a zero free domain and a chosen logarithm.  A meromorphic coefficient asymptotic is stated only for positive integral reciprocal order unless an additional branch domain is specified.

\clearpage\part{Cyclic resolution and differential transport}
\section{Cyclic projection of analytic functions}\label{sec:cyclic_projection}

\begin{readerroadmap}
The root filter asks which coefficient addresses belong to a packet. It is an exact selector. Multiplication then asks how two selected packets interact. Their residue labels add modulo $N$, so mixed inputs remain part of the product.
\end{readerroadmap}
\modelfigure{root_filter}{An exact address model for $N=3$. Every displayed exponent belongs to one residue lane. The drawing records labels rather than coefficient magnitudes.}{fig:root-filter}


Fix $N\geq1$ and write
\[
\omega=e^{2\pi i/N}.
\]
Let $\Omega$ be a disk centered at the origin, with the whole plane allowed, and assume that $\Omega$ is invariant under multiplication by $\omega$.  For a holomorphic function $F$ on $\Omega$, define
\[
(\Bcal_{N,j}F)(z)
=
\frac1N\sum_{\ell=0}^{N-1}\omega^{-j\ell}F(\omega^\ell z),
\qquad 0\leq j<N.
\]
The same formula is valid for formal power series.

\begin{theorem}[Root filter]\label{thm:root_filter}
If
\[
F(z)=\sum_{m\geq0}a_mz^m,
\]
then
\[
\Bcal_{N,j}F(z)
=
\sum_{m\equiv j\bmod N}a_mz^m.
\]
The operators satisfy
\[
\Bcal_{N,i}\Bcal_{N,j}=\delta_{ij}\Bcal_{N,j}
\]
and
\[
\sum_{j=0}^{N-1}\Bcal_{N,j}=I.
\]
\end{theorem}

\begin{proof}
Substitution of the power series gives
\[
\Bcal_{N,j}F(z)
=
\sum_{m\geq0}a_mz^m
\left(
\frac1N\sum_{\ell=0}^{N-1}\omega^{(m-j)\ell}
\right).
\]
The inner average is one when $m\equiv j\bmod N$ and zero otherwise.  The operator identities follow from the disjoint residue classes.
\end{proof}

\begin{proposition}[Covariance characterization]\label{prop:covariance}
A holomorphic function $F$ belongs to the range of $\Bcal_{N,j}$ if and only if
\[
F(\omega z)=\omega^jF(z).
\]
\end{proposition}

\begin{proof}
Every monomial selected by $\Bcal_{N,j}$ has the stated covariance.  Conversely, write $F=\sum_k\Bcal_{N,k}F$.  Replacing $z$ by $\omega z$ gives
\[
\sum_k\omega^k\Bcal_{N,k}F
=
\omega^j\sum_k\Bcal_{N,k}F.
\]
Application of $\Bcal_{N,k}$ gives
\[
(\omega^k-\omega^j)\Bcal_{N,k}F=0.
\]
The powers $1,\omega,\ldots,\omega^{N-1}$ are distinct, so
$\Bcal_{N,k}F=0$ for $k\ne j$.
\end{proof}

Multiplication is not diagonal in these coordinates. The residue labels add
modulo $N$. This is the scalar cyclic instance of the packet product from
Paper I.

\begin{theorem}[Cyclic product law]\label{thm:cyclic_product}
If $F$ and $G$ are holomorphic on the same rotation-invariant domain, then
\[
\Bcal_{N,k}(FG)
=
\sum_{\substack{0\leq i,j<N\\ i+j\equiv k\bmod N}}
(\Bcal_{N,i}F)(\Bcal_{N,j}G).
\]
\end{theorem}

\begin{proof}
Write $F=\sum_iF_i$ and $G=\sum_jG_j$, where
$F_i=\Bcal_{N,i}F$ and $G_j=\Bcal_{N,j}G$.  Their covariance gives
\[
(F_iG_j)(\omega z)=\omega^{i+j}F_i(z)G_j(z),
\]
so $F_iG_j$ belongs to the packet labelled by $i+j$ modulo $N$.
Expanding the finite product and applying $\Bcal_{N,k}$ gives the formula.
\end{proof}

\begin{remark}\label{rem:formal_analytic}
Every identity in this section holds formally.  In the holomorphic setting, the finite projector sums preserve local uniform convergence.  No choice of branch is involved.
\end{remark}

\section{The transfer exponential basis}\label{sec:transfer_exponential}

\begin{definition}\label{def:transfer_exponential}
For $\lambda\in\CC$ and $0\leq j<N$, define
\[
\Efun_{N,j}(\lambda z)
=
\Bcal_{N,j}(e^{\lambda z})
=
\sum_{n\geq0}\frac{(\lambda z)^{Nn+j}}{(Nn+j)!}.
\]
\end{definition}

\begin{theorem}[Differential cycle and initial data]\label{thm:exponential_cycle}
For $1\leq j<N$,
\[
\frac{d}{dz}\Efun_{N,j}(\lambda z)
=
\lambda\Efun_{N,j-1}(\lambda z),
\]
while
\[
\frac{d}{dz}\Efun_{N,0}(\lambda z)
=
\lambda\Efun_{N,N-1}(\lambda z).
\]
Thus
\[
\frac{d^N}{dz^N}\Efun_{N,j}(\lambda z)
=
\lambda^N\Efun_{N,j}(\lambda z).
\]
For $0\leq k<N$,
\[
\left.\frac{d^k}{dz^k}\Efun_{N,j}(\lambda z)\right|_{z=0}
=
\delta_{kj}\lambda^j.
\]
If $\lambda\ne0$, the $N$ functions $\Efun_{N,j}(\lambda z)$ form a basis of the solution space of $y^{(N)}=\lambda^Ny$.
\end{theorem}

\begin{proof}
Termwise differentiation gives the cyclic derivative formulas.  One full cycle contributes $\lambda^N$.  The first nonzero term of $\Efun_{N,j}(\lambda z)$ is $(\lambda z)^j/j!$, which gives the initial data.  The initial value vectors are linearly independent when $\lambda\ne0$.  The solution space of an order $N$ constant coefficient equation has dimension $N$, so these functions form a basis.
\end{proof}


\begin{corollary}[Wronskian of the transfer basis]\label{cor:wronskians}
Assume $N\geq2$. For $\lambda\ne0$, the Wronskian of the ordered basis
\[
\Efun_{N,0}(\lambda z),\ldots,\Efun_{N,N-1}(\lambda z)
\]
is the constant
\[
\lambda^{N(N-1)/2}.
\]
\end{corollary}

\begin{proof}
At the origin, the derivative matrix is diagonal with entries
\[
1,\lambda,\ldots,\lambda^{N-1}.
\]
For an explicit determinant proof, write each basis function as a finite
sum of exponentials with exponents $a_\ell=\lambda\omega^\ell$.
Its derivative matrix factors as
\[
W(z)=V\,\operatorname{diag}(e^{a_\ell z})\,C,
\qquad V_{k\ell}=a_\ell^k,
\]
where $C$ is the constant matrix of Fourier weights. Thus
\[
\det W(z)=e^{z\sum_\ell a_\ell}\det W(0)=\det W(0),
\]
because $N\geq2$ implies $\sum_\ell\omega^\ell=0$.
The product of the diagonal initial values is the stated power of $\lambda$.
\end{proof}

\begin{corollary}[Addition law]\label{cor:exponential_addition}
For $0\leq k<N$,
\[
\Efun_{N,k}(\lambda(z+w))
=
\sum_{\substack{0\leq i,j<N\\ i+j\equiv k\bmod N}}
\Efun_{N,i}(\lambda z)\Efun_{N,j}(\lambda w).
\]
\end{corollary}

\begin{proof}
Expand $e^{\lambda(z+w)}=e^{\lambda z}e^{\lambda w}$ as a
bivariate power series.  A term $z^a w^b$ contributes to the
$k$th section exactly when $a+b\equiv k\bmod N$.  Grouping first by
$i\equiv a\bmod N$ and $j\equiv b\bmod N$ gives the displayed sum.
\end{proof}

The basis also has a scalar hypergeometric form.

\begin{theorem}[Hypergeometric form]\label{thm:hypergeometric_form}
Let
\[
A_{N,j}
=
\left\{
\frac{j+m}{N}
\mathrel{:}
1\leq m\leq N
\text{ and }
j+m\ne N
\right\}.
\]
Then
\[
\Efun_{N,j}(\lambda z)
=
\frac{(\lambda z)^j}{j!}
{}_0F_{N-1}
\left(
\begin{matrix}-\\ A_{N,j}\end{matrix}
\mathrel{,}
\frac{(\lambda z)^N}{N^N}
\right).
\]
\end{theorem}

\begin{proof}
The multiplication formula for the gamma function gives
\[
(Nn+j)!
=
j!N^{Nn}n!
\prod_{a\in A_{N,j}}(a)_n.
\]
Indeed, the $N$ Pochhammer factors supplied by the multiplication formula
are $((j+m)/N)_n$ for $1\leq m\leq N$; the excluded factor is
$((j+N-j)/N)_n=(1)_n=n!$.
Insertion into the defining series gives the stated hypergeometric expansion.
\end{proof}

\begin{example}\label{ex:classical_exponential}
For $N=1$, the only section is $e^{\lambda z}$.  For $N=2$,
\[
\Efun_{2,0}(z)=\cosh z,
\qquad
\Efun_{2,1}(z)=\sinh z.
\]
These are classical instances of the same cyclic basis.
\end{example}

\section{Signed cyclic sections and the distinguished denominator}\label{sec:signed_system}

\begin{definition}\label{def:signed_sections}
For $0\leq j<N$, define
\[
\Dfun_{N,j}(z)
=
\sum_{n\geq0}(-1)^n\frac{z^{Nn+j}}{(Nn+j)!}.
\]
The distinguished signed denominator is
\[
\Dfun_N(z)=\Dfun_{N,0}(z).
\]
\end{definition}

Fix any $\rho\in\CC$ with $\rho^N=-1$.

\begin{proposition}[Relation with the transfer exponential]\label{prop:signed_exponential_relation}
For every $j$,
\[
\Dfun_{N,j}(z)
=
\rho^{-j}\Efun_{N,j}(\rho z).
\]
\end{proposition}

\begin{proof}
The coefficient of $z^{Nn+j}$ on the right is
\[
\rho^{-j}\rho^{Nn+j}\frac1{(Nn+j)!}
=
(-1)^n\frac1{(Nn+j)!}.
\]
\end{proof}

\begin{theorem}[Signed derivative cycle]\label{thm:signed_cycle}
For $1\leq j<N$,
\[
\Dfun_{N,j}'=\Dfun_{N,j-1},
\]
and
\[
\Dfun_{N,0}'=-\Dfun_{N,N-1}.
\]
Every signed section satisfies
\[
\Dfun_{N,j}^{(N)}=-\Dfun_{N,j}.
\]
The $N$ signed sections form the basis of $y^{(N)}=-y$ determined by
\[
\Dfun_{N,j}^{(k)}(0)=\delta_{jk}
\]
for $0\leq j,k<N$.
\end{theorem}

\begin{proof}
The derivative formulas follow directly from the series.  The first nonzero term of $\Dfun_{N,j}$ is $z^j/j!$, so the initial value matrix is the identity.  The basis statement follows from the dimension of the solution space.
\end{proof}

\begin{corollary}[Wronskian of the signed basis]\label{cor:signed_wronskian}
Assume $N\geq2$. The Wronskian of
\[
\Dfun_{N,0}(z),\ldots,\Dfun_{N,N-1}(z)
\]
is one.
\end{corollary}

\begin{proof}
The derivative matrix at the origin is the identity. The determinant
factorization in the proof of \cref{cor:wronskians} applies with
$a_\ell=\rho\omega^\ell$ and with the signed-section Fourier weights.
Since $\sum_\ell a_\ell=0$, the determinant is constant and equals one.
\end{proof}


\begin{proposition}[Twisted addition law]\label{prop:signed_addition}
For $0\leq k<N$,
\[
\Dfun_{N,k}(z+w)
=
\sum_{\substack{0\leq i,j<N\\ i+j\equiv k\bmod N}}
(-1)^{\lfloor(i+j)/N\rfloor}
\Dfun_{N,i}(z)\Dfun_{N,j}(w).
\]
\end{proposition}

\begin{proof}
Use \cref{prop:signed_exponential_relation,cor:exponential_addition}.  The factor introduced by a pair $(i,j)$ is
\[
\rho^{i+j-k}.
\]
Since $0\leq i+j\leq2N-2$ and $i+j\equiv k\bmod N$, this factor is one when $i+j=k$ and $-1$ when $i+j=k+N$.
\end{proof}

\begin{theorem}[Finite exponential representation]\label{thm:finite_exponential_representation}
Let
\[
\Rcal_N=\{\eta\in\CC\mathrel{:}\eta^N=-1\}.
\]
Then
\[
\Dfun_N(z)
=
\frac1N\sum_{\eta\in\Rcal_N}e^{\eta z}.
\]
More generally,
\[
\Dfun_{N,j}(z)
=
\frac{\rho^{-j}}N
\sum_{\ell=0}^{N-1}
\omega^{-j\ell}e^{\rho\omega^\ell z}.
\]
\end{theorem}

\begin{proof}
The second formula is the root filter expression for $\rho^{-j}\Efun_{N,j}(\rho z)$.  The first is the case $j=0$.
\end{proof}

\begin{definition}\label{def:reduced_denominator}
Define the reduced denominator
\[
\Delta_N(w)
=
\sum_{n\geq0}(-1)^n\frac{w^n}{(Nn)!}.
\]
Then
\[
\Dfun_N(z)=\Delta_N(z^N).
\]
\end{definition}

\begin{corollary}[Symmetries]\label{cor:denominator_symmetries}
The denominator satisfies
\[
\Dfun_N(\omega z)=\Dfun_N(z)
\]
and
\[
\Dfun_N(\overline z)=\overline{\Dfun_N(z)}.
\]
Its nonzero zeros occur in orbits under multiplication by the $N$th roots of unity and under complex conjugation.
\end{corollary}

\begin{proof}
The first identity follows from the powers $z^{Nn}$.  The second follows from the real coefficients.  Both symmetries preserve the zero set and multiplicity.
\end{proof}

\begin{example}\label{ex:small_denominators}
The first denominators are
\[
\Dfun_1(z)=e^{-z},
\qquad
\Dfun_2(z)=\cos z.
\]
For $N=3$,
\[
\Dfun_3(z)
=
\frac13\left(
e^{-z}
+2e^{z/2}\cos\left(\frac{\sqrt3z}{2}\right)
\right).
\]
For $N=4$,
\[
\Dfun_4(z)
=
\cos\left(\frac z{\sqrt2}\right)
\cosh\left(\frac z{\sqrt2}\right).
\]
\end{example}

\section{A finite algebra and matrix model}\label{sec:finite_algebra}

\begin{keyidea}
The signed shift is a finite carrier for the differential cycle. At order four, three basis steps have coefficient one and the final step has coefficient minus one. Four steps return the vector with reversed sign. The matrix exponential packages every finite path length in one analytic expression.
\end{keyidea}
\modelfigure{signed_cycle}{The exact signed shift $J_4$. Source columns and target rows describe the same action as the cycle on the left. Direct integer multiplication verifies $J_4^4=-I$.}{fig:signed-cycle}


The derivative and addition laws can be read from one finite algebra identity.  This formulation also identifies the exponent polygon as the spectrum of a signed cyclic shift.

\begin{definition}\label{def:cyclic_algebras}
Let
\[
\mathfrak A_N^+=\CC[\xi]/(\xi^N-1)
\]
and
\[
\mathfrak A_N^-=\CC[\xi]/(\xi^N+1).
\]
Every element has a unique representative of degree below $N$.
\end{definition}

\begin{theorem}[Finite algebra exponential identities]\label{thm:finite_algebra_exponentials}
In $\mathfrak A_N^+$ one has
\[
e^{\lambda z\xi}
=
\sum_{j=0}^{N-1}\Efun_{N,j}(\lambda z)\xi^j.
\]
In $\mathfrak A_N^-$ one has
\[
e^{z\xi}
=
\sum_{j=0}^{N-1}\Dfun_{N,j}(z)\xi^j.
\]
Differentiation and multiplication of these identities recover the cyclic derivative systems and the two addition laws.
\end{theorem}

\begin{proof}
Expand the exponential and write every exponent uniquely as $Nn+j$.  In $\mathfrak A_N^+$, the relation $\xi^{Nn+j}=\xi^j$ gives the transfer sections.  In $\mathfrak A_N^-$, the relation $\xi^{Nn+j}=(-1)^n\xi^j$ gives the signed sections.

Differentiation multiplies by $\lambda\xi$ in the first algebra and by $\xi$ in the second.  Reduction of the product by the defining relation gives the derivative cycles.  The identity
\[
e^{(z+w)\xi}=e^{z\xi}e^{w\xi}
\]
gives the addition laws after comparison of the coefficients of $1,\xi,\ldots,\xi^{N-1}$.
\end{proof}

\begin{definition}\label{def:signed_shift_matrix}
Let $J_N$ be the signed cyclic shift on the standard basis $e_0,\ldots,e_{N-1}$ of $\CC^N$, defined by
\[
J_Ne_j=e_{j+1}
\quad (0\leq j<N-1),
\qquad
J_Ne_{N-1}=-e_0.
\]
Then $J_N^N=-I$.
\end{definition}

\begin{theorem}[Companion matrix realization]\label{thm:companion_matrix_realization}
Put
\[
\mathbf D_N(z)
=
\begin{pmatrix}
\Dfun_{N,0}(z)\\
\Dfun_{N,1}(z)\\
\vdots\\
\Dfun_{N,N-1}(z)
\end{pmatrix}.
\]
Then
\[
\mathbf D_N(z)=e^{zJ_N}e_0
\]
and
\[
\Dfun_N(z)=\frac1N\tr(e^{zJ_N}).
\]
The eigenvalues of $J_N$ are the $N$ roots of $-1$, each with multiplicity one.
\end{theorem}

\begin{proof}
The vector $e^{zJ_N}e_0$ has derivative $J_Ne^{zJ_N}e_0$ and initial value $e_0$.  Its components therefore satisfy the same signed derivative cycle and initial data as $\mathbf D_N$.  Uniqueness for the linear system gives the first identity.

The vectors
\[
e_0,J_Ne_0,\ldots,J_N^{N-1}e_0
\]
are the standard basis.  Thus $e_0$ is cyclic for $J_N$, so the
minimal polynomial of $J_N$ has degree $N$.  Since $J_N^N=-I$, that
minimal polynomial divides $x^N+1$ and therefore equals $x^N+1$.
The characteristic polynomial has the same degree and is divisible
by the minimal polynomial, so it too equals $x^N+1$.  This polynomial
has distinct roots.  Consequently $J_N$ is diagonalizable and its
spectrum consists of all $N$ roots of $-1$, each once.  Hence
\[
\tr(e^{zJ_N})
=
\sum_{\eta^N=-1}e^{\eta z}
=
N\Dfun_N(z)
\]
by \cref{thm:finite_exponential_representation}.
\end{proof}

\begin{corollary}[Entries of the fundamental matrix]\label{cor:fundamental_matrix_entries}
For $0\leq p,q<N$, let $d_{pq}$ be the member of $\{0,\ldots,N-1\}$ congruent to $p-q$ modulo $N$.  Put
\[
\varepsilon_{pq}
=
\begin{cases}
1,&p\geq q,\\
-1,&p<q.
\end{cases}
\]
Then
\[
\bigl(e^{zJ_N}\bigr)_{pq}
=
\varepsilon_{pq}\Dfun_{N,d_{pq}}(z).
\]
Thus every entry of the fundamental matrix is one of the signed sections, with the sign recording whether the cyclic path from $q$ to $p$ crosses the signed edge.
\end{corollary}

\begin{proof}
For every nonnegative integer $n$,
\[
J_N^{Nn+d_{pq}}e_q
=
(-1)^n\varepsilon_{pq}e_p.
\]
Only powers whose exponent is congruent to $d_{pq}$ modulo $N$ contribute to the $(p,q)$ entry of the exponential.  Therefore
\[
\bigl(e^{zJ_N}\bigr)_{pq}
=
\varepsilon_{pq}
\sum_{n\geq0}(-1)^n
\frac{z^{Nn+d_{pq}}}{(Nn+d_{pq})!}.
\]
This is the stated formula.
\end{proof}

\begin{corollary}[Addition and inversion laws for the companion matrix]\label{cor:companion_addition_inversion}
Assume $N\geq2$.  The matrix
\[
U_N(z)=e^{zJ_N}
\]
satisfies
\[
U_N(z+w)=U_N(z)U_N(w),
\qquad
U_N(z)^{-1}=U_N(-z),
\qquad
\det U_N(z)=1.
\]
Its first column is $\mathbf D_N(z)$, and all of its entries are given by \cref{cor:fundamental_matrix_entries}.
\end{corollary}

\begin{proof}
The first two identities are the exponential group law.  Since $\tr J_N=0$,
\[
\det(e^{zJ_N})=e^{z\tr J_N}=1.
\]
The first column statement is \cref{thm:companion_matrix_realization}.
\end{proof}

\clearpage\part{Reciprocal coefficients and polynomial families}
\section{Reciprocal powers and parameter polynomials}\label{sec:reciprocal_powers}

\begin{readerroadmap}
A unit constant term makes the formal reciprocal unambiguous. To compute one coefficient, only finitely many earlier coefficients are needed. The later analytic theory explains how those exact coefficients relate to the full zero set.
\end{readerroadmap}
\modelfigure{reciprocal_normalization}{At $N=2$ and $r=1$, the reciprocal is $\sec z$. The upper row gives block exponential coefficients and the lower row gives ordinary Taylor coefficients. Both rows encode the same five terms.}{fig:reciprocal-normalization}


Since $\Dfun_N(0)=1$, its formal logarithm is defined at the origin.  On a
simply connected zero free region containing the origin, there is a unique
holomorphic branch normalized by $\log\Dfun_N(0)=0$, and its germ at the
origin agrees with the formal logarithm.

\begin{definition}\label{def:reciprocal_family}
For $r\in\CC$, define
\[
\Sfun_N^{(r)}(z)
=
\Dfun_N(z)^{-r}
=
\exp\bigl(-r\log\Dfun_N(z)\bigr).
\]
Write
\[
\Sfun_N^{(r)}(z)
=
\sum_{n\geq0}P_{N,n}(r)\frac{z^{Nn}}{(Nn)!}.
\]
Define the logarithmic coefficients $\kappa_{N,m}$ by
\[
-\log\Delta_N(w)
=
\sum_{m\geq1}\kappa_{N,m}\frac{w^m}{m}.
\]
\end{definition}

The polynomial dependence on $r$ can be read directly from the formal logarithm.  We use the complete exponential Bell polynomials $B_n$ defined by
\[
\exp\left(\sum_{m\geq1}x_m\frac{t^m}{m!}\right)
=
\sum_{n\geq0}B_n(x_1,\ldots,x_n)\frac{t^n}{n!}.
\]

\begin{theorem}[Bell polynomial formula]\label{thm:reciprocal_bell}
For every $n\geq0$,
\[
P_{N,n}(r)
=
\frac{(Nn)!}{n!}
B_n\bigl(
 r\kappa_{N,1},
 r\,1!\,\kappa_{N,2},
 \ldots,
 r\,(n-1)!\,\kappa_{N,n}
\bigr).
\]
For $n\geq1$, this is a polynomial in $r$ of exact degree $n$.  Its leading coefficient is
\[
\frac{(Nn)!}{n!(N!)^n}.
\]
\end{theorem}

\begin{proof}
The identity
\[
\Delta_N(w)^{-r}
=
\exp\left(
 r\sum_{m\geq1}\kappa_{N,m}\frac{w^m}{m}
\right)
\]
fits the defining generating function for $B_n$ with
\[
x_m=r(m-1)!\kappa_{N,m}.
\]
The coefficient of $w^n$ is therefore the displayed Bell polynomial divided by $n!$.  Multiplication by $(Nn)!$ gives $P_{N,n}(r)$.  Since
\[
\Delta_N(w)=1-\frac{w}{N!}+O(w^2),
\]
we have $\kappa_{N,1}=1/N!$.  The term $x_1^n$ gives the stated leading coefficient, which is nonzero.
\end{proof}

The coefficients also satisfy a triangular recurrence that does not require the logarithmic coefficients to be known first.

\begin{theorem}[Parameter recurrence]\label{thm:parameter_recurrence}
We have $P_{N,0}(r)=1$.  For $n\geq1$,
\[
\begin{aligned}
P_{N,n}(r)
={}&-
\sum_{k=1}^{n-1}
(-1)^k
\binom{Nn-1}{Nk}
P_{N,n-k}(r)
\\
&-r
\sum_{k=1}^{n}
(-1)^k
\binom{Nn-1}{Nk-1}
P_{N,n-k}(r).
\end{aligned}
\]
\end{theorem}

\begin{proof}
Set $S=\Dfun_N^{-r}$.  Logarithmic differentiation gives
\[
\Dfun_NS'+r\Dfun_N'S=0.
\]
Taking the coefficient of $z^{Nn-1}/(Nn-1)!$ gives
\[
\sum_{k=0}^{n-1}(-1)^k
\binom{Nn-1}{Nk}P_{N,n-k}(r)
+r\sum_{k=1}^{n}(-1)^k
\binom{Nn-1}{Nk-1}P_{N,n-k}(r)=0.
\]
The term with $k=0$ in the first sum is $P_{N,n}(r)$, and every
other term has a smaller block index.  Rearrangement gives the formula.
\end{proof}

\begin{theorem}[Additive parameter law for reciprocal powers]\label{thm:reciprocal_parameter_addition}
For all $r,s\in\CC$,
\[
\Sfun_N^{(r)}(z)\Sfun_N^{(s)}(z)
=
\Sfun_N^{(r+s)}(z)
\]
as a formal identity and on every common zero-free region with a fixed logarithm.  Consequently,
\[
P_{N,n}(r+s)
=
\sum_{q=0}^{n}
\binom{Nn}{Nq}
P_{N,q}(r)P_{N,n-q}(s).
\]
\end{theorem}

\begin{proof}
Both statements follow by adding the exponents of the same logarithm and then multiplying the block exponential generating series.
\end{proof}

\begin{corollary}[A sequence of binomial type]\label{cor:binomial_type}
Define
\[
A_{N,n}(r)
=
\frac{n!}{(Nn)!}P_{N,n}(r).
\]
Then
\[
A_{N,n}(r+s)
=
\sum_{q=0}^{n}
\binom nq
A_{N,q}(r)A_{N,n-q}(s).
\]
Thus $(A_{N,n})_{n\geq0}$ is a polynomial sequence of binomial type.
\end{corollary}

\begin{proof}
Multiply the coefficient identity in \cref{thm:reciprocal_parameter_addition} by $n!/(Nn)!$ and simplify the factorials.
\end{proof}


The recurrence in \cref{thm:parameter_recurrence} uses the coefficients of the original denominator.  A second recurrence uses only the logarithmic coefficients and is often shorter.

\begin{theorem}[Logarithmic Newton recurrence]\label{thm:logarithmic_newton_recurrence}
Set
\[
b_{N,n}(r)=\frac{P_{N,n}(r)}{(Nn)!}.
\]
Then $b_{N,0}(r)=1$ and
for every $n\geq1$,
\[
nb_{N,n}(r)
=
r\sum_{k=1}^{n}\kappa_{N,k}b_{N,n-k}(r).
\]
Equivalently,
\[
P_{N,n}(r)
=
\frac r n
\sum_{k=1}^{n}
\frac{(Nn)!}{(N(n-k))!}
\kappa_{N,k}P_{N,n-k}(r).
\]
If $p_k=r\kappa_{N,k}$, then, for $n\geq1$,
\[
b_{N,n}(r)
=
\frac1{n!}
\det
\begin{pmatrix}
p_1&-1&0&\cdots&0\\
p_2&p_1&-2&\cdots&0\\
p_3&p_2&p_1&\ddots&\vdots\\
\vdots&\vdots&\ddots&\ddots&-(n-1)\\
p_n&p_{n-1}&\cdots&p_2&p_1
\end{pmatrix}.
\]
\end{theorem}

\begin{proof}
Let
\[
B_r(w)=\Delta_N(w)^{-r}
=
\sum_{n\geq0}b_{N,n}(r)w^n.
\]
Logarithmic differentiation gives
\[
\frac{B_r'(w)}{B_r(w)}
=
r\sum_{k\geq1}\kappa_{N,k}w^{k-1}.
\]
Comparison of the coefficient of $w^{n-1}$ proves the recurrence.

Let $D_n$ denote the displayed determinant, with $D_0=1$.  Expansion along the last row gives
\[
D_n
=
\sum_{k=1}^{n}
\frac{(n-1)!}{(n-k)!}p_kD_{n-k}.
\]
Here the cofactor sign is cancelled by the product of the negative
superdiagonal entries $-(n-k+1),\ldots,-(n-1)$.
Thus $D_n/n!$ satisfies the same recurrence and initial value as $b_{N,n}(r)$.
\end{proof}

\begin{corollary}[The first three reciprocal polynomials]\label{cor:first_reciprocal_polynomials}
For every $N\geq1$,
\[
P_{N,1}(r)=r,
\]
\[
P_{N,2}(r)
=
r\left[
\frac12\binom{2N}{N}(r+1)-1
\right],
\]
and
\[
P_{N,3}(r)
=
r\left[
\frac{(3N)!}{6(N!)^3}(r+1)(r+2)
-
\binom{3N}{N}(r+1)
+1
\right].
\]
\end{corollary}

\begin{proof}
The first logarithmic coefficients are
\[
\kappa_{N,1}=\frac1{N!},
\]
\[
\kappa_{N,2}=\frac1{(N!)^2}-\frac2{(2N)!},
\]
and
\[
\kappa_{N,3}
=
\frac1{(N!)^3}
-
\frac3{N!(2N)!}
+
\frac3{(3N)!}.
\]
Insert these values into \cref{thm:reciprocal_bell} and simplify.
\end{proof}

\begin{corollary}\label{cor:factor_r}
For $n\geq1$, the polynomial $P_{N,n}(r)$ is divisible by $r$.
\end{corollary}

\begin{proof}
At $r=0$, the reciprocal family is the constant function one.  Every positive-index block coefficient therefore vanishes.
\end{proof}

\begin{example}\label{ex:reciprocal_polynomials}
For $N=1$, one has $\Dfun_1(z)^{-r}=e^{rz}$ and $P_{1,n}(r)=r^n$.  For $N=2$,
\[
P_{2,1}(r)=r,
\]
\[
P_{2,2}(r)=r(3r+2),
\]
and
\[
P_{2,3}(r)=r(15r^2+30r+16).
\]
At $r=1$, these are the positive secant coefficients.
\end{example}

\section{Logarithmic derivatives and Riccati equations}\label{sec:tangent}

\begin{microexample}[title={A coefficient calculation at order two}]
Write $\Delta_2(w)^{-r}=\sum_{n\geq0}b_n(r)w^n$. Since
\[
\Delta_2(w)=1-\frac w2+\frac{w^2}{24}+O(w^3),
\qquad -\log\Delta_2(w)=\frac w2+\frac{w^2}{12}+O(w^3),
\]
the first logarithmic coefficients are $\kappa_{2,1}=1/2$ and $\kappa_{2,2}=1/6$. Newton's recurrence gives
\[
b_1(r)=\frac r2,\qquad
2b_2(r)=r\left(\frac r4+\frac16\right).
\]
Restoring the factorial normalization yields $P_{2,1}(r)=r$ and $P_{2,2}(r)=r(3r+2)$. At $r=1$, the second value is five. No zero-location assumption is needed for this finite calculation.
\end{microexample}
\modelfigure{parameter_polynomials}{The two exact parameter polynomials from the worked calculation. Each plotted curve evaluates its displayed formula. The points at $r=1$ recover the first two positive-index secant coefficients.}{fig:parameter-polynomials}


\begin{definition}\label{def:transfer_tangent}
Define the transfer tangent
\[
\Tfun_N(z)
=-\frac{\Dfun_N'(z)}{\Dfun_N(z)}.
\]
Near the origin, its convergent block expansion has the form
\[
\Tfun_N(z)
=
\sum_{n\geq0}T_N(n)
\frac{z^{Nn+N-1}}{(Nn+N-1)!}.
\]
\end{definition}

\begin{proposition}[Logarithmic coefficient formula]\label{prop:tangent_kappa}
The transfer tangent satisfies
\[
\Tfun_N(z)
=
N\sum_{m\geq1}\kappa_{N,m}z^{Nm-1}.
\]
Equivalently,
\[
T_N(n)
=
N(Nn+N-1)!\kappa_{N,n+1}.
\]
\end{proposition}

\begin{proof}
Since $\Dfun_N(z)=\Delta_N(z^N)$,
\[
-\frac{\Dfun_N'(z)}{\Dfun_N(z)}
=
-Nz^{N-1}\frac{\Delta_N'(z^N)}{\Delta_N(z^N)}.
\]
Differentiate the defining series for $-\log\Delta_N$ and compare coefficients.
\end{proof}

\begin{theorem}[Tangent recurrence]\label{thm:tangent_recurrence}
The coefficients satisfy $T_N(0)=1$ and
\[
T_N(n)
=
(-1)^n
-
\sum_{k=1}^{n}
(-1)^k
\binom{Nn+N-1}{Nk}
T_N(n-k).
\]
\end{theorem}

\begin{proof}
The defining equation is
\[
\Dfun_N\Tfun_N=-\Dfun_N'.
\]
Extracting the coefficient of $z^{Nn+N-1}/(Nn+N-1)!$ gives
\[
\sum_{k=0}^{n}(-1)^k
\binom{Nn+N-1}{Nk}T_N(n-k)=(-1)^n.
\]
Isolating the $k=0$ term gives the recurrence.
\end{proof}

\begin{corollary}[Derivative of the reciprocal family]\label{cor:reciprocal_derivative}
For every $r\in\CC$, the formal identity
\[
\frac{d}{dz}\Sfun_N^{(r)}(z)
=
r\Tfun_N(z)\Sfun_N^{(r)}(z)
\]
holds.  The same identity holds analytically on every zero-free
domain on which the logarithm used to define $\Sfun_N^{(r)}$ is fixed.
\end{corollary}

\begin{proof}
Differentiate $\Dfun_N^{-r}$ on a zero-free region.  The formal identity follows from the same logarithmic calculation.
\end{proof}

For a function $g$, define the complete differential Bell polynomial $Y_N$ by
\[
\frac{d^N}{dz^N}e^{g(z)}
=
e^{g(z)}Y_N\bigl(g',g'',\ldots,g^{(N)}\bigr).
\]

\begin{theorem}[Scalar higher Riccati identity]\label{thm:higher_riccati}
At every point of the open set $\{z:\Dfun_N(z)\ne0\}$,
\[
Y_N\bigl(
-\Tfun_N,
-\Tfun_N',
\ldots,
-\Tfun_N^{(N-1)}
\bigr)
=-1.
\]
\end{theorem}

\begin{proof}
Write $\Dfun_N=e^g$ locally.  Then $g'=-\Tfun_N$.  Divide the equation $\Dfun_N^{(N)}=-\Dfun_N$ by $\Dfun_N$ and use the defining identity for $Y_N$.
\end{proof}

This scalar identity follows from one ordinary differential equation. It does
not redevelop the general nonlinear operation-tree theory of Paper III.

\begin{example}\label{ex:ordinary_tangent}
For $N=1$, $\Tfun_1=1$.  For $N=2$, $\Dfun_2=\cos z$ and $\Tfun_2=\tan z$.  The Riccati equation becomes
\[
\Tfun_2'=1+\Tfun_2^2.
\]
\end{example}

\section{Residue-section reciprocals and Appell systems}\label{sec:residue_section}

The section $\Efun_{N,j}$ has a forced zero of order $j$ at the origin.  Removing it leaves a unit series in $z^N$.

\begin{definition}\label{def:normalized_section}
For $0\leq j<N$, define
\[
\Phi_{N,j}(w)
=
j!\sum_{n\geq0}\frac{w^n}{(Nn+j)!}.
\]
Then
\[
\Efun_{N,j}(z)
=
\frac{z^j}{j!}\Phi_{N,j}(z^N).
\]
Define the residue-section reciprocal function
\[
\Cfun_{N,j}(z)
=
\Phi_{N,j}(z^N)^{-1}
=
\frac{z^j/j!}{\Efun_{N,j}(z)}.
\]
For $j>0$, the last quotient at $z=0$ is interpreted by its removable
continuation, namely the first expression in the display.
Write
\[
\Cfun_{N,j}(z)
=
\sum_{n\geq0}C_{N,j}(n)\frac{z^{Nn}}{(Nn)!}.
\]
\end{definition}

\begin{theorem}[Triangular inverse system]\label{thm:triangular_inverse}
The coefficients are determined by $C_{N,j}(0)=1$ and
\[
\sum_{q=0}^{n}
\binom{Nn+j}{Nq}
C_{N,j}(q)
=
\delta_{n0}.
\]
For $n\geq1$,
\[
C_{N,j}(n)
=
-
\frac{
\displaystyle\sum_{q=0}^{n-1}
\binom{Nn+j}{Nq}C_{N,j}(q)
}{
\displaystyle\binom{Nn+j}{Nn}
}.
\]
\end{theorem}

\begin{proof}
Multiply the defining series for $\Cfun_{N,j}$ and $\Efun_{N,j}$.  Their product is $z^j/j!$.  At degree $Nn+j$, multiplication by $(Nn+j)!$ gives the stated binomial coefficient.  The diagonal coefficient is nonzero, so the system is triangular.
\end{proof}

\begin{definition}\label{def:residue_section_power}
For $r\in\CC$, define formal coefficients $C_{N,j,n}(r)$ by
\[
\Cfun_{N,j}(z)^r
=
\sum_{n\geq0}C_{N,j,n}(r)\frac{z^{Nn}}{(Nn)!}.
\]
On a simply connected region containing $0$ and avoiding the zeros of $\Phi_{N,j}(z^N)$, interpret this power as $e^{rL}$, where $L$ is the holomorphic logarithm of $\Cfun_{N,j}$ normalized by $L(0)=0$.
\end{definition}

\begin{theorem}[Parameter-addition law for residue-section reciprocals]\label{thm:residue_section_parameter_addition}
For each $n\geq1$, the coefficient $C_{N,j,n}(r)$ is a polynomial in $r$ of exact degree $n$.  Its leading coefficient is
\[
\frac{(Nn)!}{n!}
\left(-\frac{j!}{(N+j)!}\right)^n.
\]
It satisfies
\[
C_{N,j,n}(r+s)
=
\sum_{q=0}^{n}
\binom{Nn}{Nq}
C_{N,j,q}(r)C_{N,j,n-q}(s).
\]
In particular,
\[
C_{N,j,n}(1)=C_{N,j}(n).
\]
\end{theorem}

\begin{proof}
The unit series $\Phi_{N,j}(w)^{-1}$ has a formal logarithm with zero constant term.  Its linear term is
\[
-\frac{j!}{(N+j)!}w.
\]
Exponentiation after multiplication by $r$ shows that the coefficient of $w^n$ is a polynomial of degree at most $n$.  The $n$th power of the linear logarithmic term gives the displayed nonzero leading coefficient, so the degree is exactly $n$.  Multiplication of the powers with parameters $r$ and $s$ gives the coefficient identity.
More explicitly, before the factorial normalization the coefficient of
$r^nw^n$ is
\[
\frac1{n!}\left(-\frac{j!}{(N+j)!}\right)^n;
\]
thus multiplication by $(Nn)!$ gives the displayed leading coefficient.
\end{proof}

\begin{definition}\label{def:block_appell}
Define polynomials $E_m^{(N,j,r)}(x)$ by
\[
\Cfun_{N,j}(z)^re^{xz}
=
\sum_{m\geq0}E_m^{(N,j,r)}(x)\frac{z^m}{m!}.
\]
\end{definition}

\begin{theorem}[Appell laws]\label{thm:appell_laws}
For $m\geq1$, the polynomials satisfy
\[
\frac{d}{dx}E_m^{(N,j,r)}(x)
=
mE_{m-1}^{(N,j,r)}(x)
\]
and, for $m\geq0$,
\[
E_m^{(N,j,r)}(x+y)
=
\sum_{q=0}^{m}
\binom mq
E_q^{(N,j,r)}(x)y^{m-q}.
\]
\end{theorem}

\begin{proof}
Differentiate the generating function with respect to $x$ for the first identity.  Multiply it by $e^{yz}$ and compare coefficients for the second.
\end{proof}

\begin{proposition}[Explicit Appell polynomial]\label{prop:explicit_appell}
For every $m\geq0$,
\[
E_m^{(N,j,r)}(x)
=
\sum_{0\leq n\leq m/N}
\binom{m}{Nn}
C_{N,j,n}(r)x^{m-Nn}.
\]
\end{proposition}

\begin{proof}
Multiply the block series for $\Cfun_{N,j}^r$ by the exponential series and collect the coefficient of total degree $m$.
\end{proof}


\begin{theorem}[Joint Appell convolution in the space and order parameters]\label{thm:appell_two_parameter_convolution}
For all $r,s,x,y\in\CC$ and $m\geq0$,
\[
E_m^{(N,j,r+s)}(x+y)
=
\sum_{q=0}^{m}
\binom mq
E_q^{(N,j,r)}(x)
E_{m-q}^{(N,j,s)}(y).
\]
\end{theorem}

\begin{proof}
Multiply the generating functions
\[
\Cfun_{N,j}(z)^re^{xz}
\quad\hbox{and}\quad
\Cfun_{N,j}(z)^se^{yz}.
\]
Their product is the generating function with parameters $r+s$ and $x+y$.  Coefficient comparison gives the formula.
\end{proof}

\begin{proposition}[Formal differential-operator representation on polynomials]\label{prop:appell_operator_form}
Let a formal series $A(z)=\sum_{q\geq0}a_qz^q/q!$ act on polynomials by
\[
A(\partial_x)f(x)
=
\sum_{q\geq0}\frac{a_q}{q!}\partial_x^qf(x).
\]
The power $r$ is applied to the formal symbol before the substitution $z=\partial_x$.  Then
\[
E_m^{(N,j,r)}(x)
=
\Cfun_{N,j}(\partial_x)^r x^m.
\]
The inverse relation is
\[
x^m
=
\Phi_{N,j}(\partial_x^N)^r
E_m^{(N,j,r)}(x).
\]
Both sums are finite on every polynomial.
\end{proposition}

\begin{proof}
Applying $A(\partial_x)$ to $e^{xz}$ multiplies it by $A(z)$.  Taking
$A=\Cfun_{N,j}^r$ gives the first identity by the defining generating
function.  Constant coefficient operators on polynomials satisfy
\[
(AB)(\partial_x)=A(\partial_x)B(\partial_x),
\]
because all derivative sums are finite and powers of $\partial_x$
commute.  The two formal symbols
\[
\Cfun_{N,j}(z)^r=\Phi_{N,j}(z^N)^{-r}
\quad\hbox{and}\quad
\Phi_{N,j}(z^N)^r
\]
multiply to one.  Their substitutions at $z=\partial_x$ are therefore
inverse constant coefficient operators on polynomials, which proves
the second identity.
\end{proof}


\begin{corollary}[Appell symmetry]\label{cor:appell_symmetry}
Each polynomial $E_m^{(N,j,r)}$ is monic of degree $m$, and
\[
E_m^{(N,j,r)}(\omega x)
=
\omega^mE_m^{(N,j,r)}(x).
\]
\end{corollary}

\begin{proof}
The term $n=0$ in \cref{prop:explicit_appell} is $x^m$.  Every other exponent is $m-Nn$, so replacement of $x$ by $\omega x$ multiplies every term by $\omega^m$.
\end{proof}

\begin{example}\label{ex:euler_boundaries}
For $(N,j)=(2,0)$,
\[
\Cfun_{2,0}(z)=\sech z.
\]
For $(N,j)=(2,1)$,
\[
\Cfun_{2,1}(z)=\frac z{\sinh z}.
\]
The first case is a scaled and translated Euler family.  In fact, if $E_m(x)$ denotes the classical Euler polynomial, then
\[
E_m^{(2,0,1)}(x)
=
2^mE_m\left(\frac{x+1}{2}\right).
\]
If $B_m(x)$ denotes the classical Bernoulli polynomial, the second case has
the exact identification
\[
E_m^{(2,1,1)}(x)
=
2^mB_m\left(\frac{x+1}{2}\right),
\]
obtained by substituting $t=2z$ in
$t e^{yt}/(e^t-1)$.
\end{example}

\section{Scaled reciprocal coupling and theta polynomials}\label{sec:theta}

\begin{definition}\label{def:unsigned_denominator}
Set
\[
H_N(z)=\Efun_{N,0}(z)
=
\sum_{n\geq0}\frac{z^{Nn}}{(Nn)!}.
\]
For $u\in\CC$, define $\Thetafun_{N,n}(u)$ by
\[
H_N(uz)^{-1}H_N(z)^{-u}
=
\sum_{n\geq0}\Thetafun_{N,n}(u)\frac{z^{Nn}}{(Nn)!}.
\]
This is first a formal identity at the origin.  Analytically, it is interpreted
on a common simply connected zero-free region with the branch of
$\log H_N$ normalized by $\log H_N(0)=0$.
\end{definition}

\begin{theorem}[Theta reconstruction]\label{thm:theta_reconstruction}
For every $n\geq0$,
\[
\Thetafun_{N,n}(u)
=
(-1)^n
\sum_{q=0}^{n}
\binom{Nn}{Nq}
P_{N,q}(1)u^{Nq}P_{N,n-q}(u).
\]
Thus $\Thetafun_{N,n}$ is a polynomial of degree at most $Nn$.
\end{theorem}

\begin{proof}
Choose $\rho$ with $\rho^N=-1$.  Then
\[
H_N(z)=\Dfun_N(\rho z).
\]
The coefficient of $z^{Nn}$ in $H_N(z)^{-r}$ is therefore $(-1)^nP_{N,n}(r)/(Nn)!$.  Apply this to $H_N(uz)^{-1}$ and $H_N(z)^{-u}$, then multiply the two block series.
\end{proof}

\begin{proposition}[Trivial parameter factors]\label{prop:theta_factors}
For $n\geq1$, the polynomial $\Thetafun_{N,n}(u)$ is divisible by $u$.  If $N$ is even, it is also divisible by $u+1$.
\end{proposition}

\begin{proof}
At $u=0$, the defining product is one.  Hence every positive-index coefficient vanishes.  If $N$ is even, then $H_N(-z)=H_N(z)$.  At $u=-1$, the product is again one, so every positive-index coefficient vanishes there as well.
\end{proof}


\begin{proposition}[Theta Newton recurrence]\label{prop:theta_newton_recurrence}
Put
\[
\theta_{N,n}(u)=\frac{\Thetafun_{N,n}(u)}{(Nn)!}.
\]
Then $\theta_{N,0}(u)=1$ and
\[
n\theta_{N,n}(u)
=
\sum_{m=1}^{n}
(-1)^m\kappa_{N,m}
\bigl(u^{Nm}+u\bigr)
\theta_{N,n-m}(u).
\]
In particular,
\[
\Thetafun_{N,1}(u)=-(u^N+u).
\]
\end{proposition}

\begin{proof}
Write $w=z^N$.  Since $H_N(z)=\Delta_N(-w)$,
\[
-\log H_N(z)
=
\sum_{m\geq1}(-1)^m\kappa_{N,m}\frac{w^m}{m}.
\]
The logarithm of the theta generating function is therefore
\[
\sum_{m\geq1}
(-1)^m\kappa_{N,m}
\bigl(u^{Nm}+u\bigr)
\frac{w^m}{m}.
\]
Logarithmic differentiation and coefficient comparison give the recurrence.  The case $n=1$ uses $\kappa_{N,1}=1/N!$.
\end{proof}

\begin{example}\label{ex:theta_even}
For $N=2$,
\[
\Thetafun_{2,1}(u)=-u(u+1)
\]
and
\[
\Thetafun_{2,2}(u)
=
u(u+1)(5u^2+u+2).
\]
The two factors predicted by \cref{prop:theta_factors} are visible.
\end{example}

\clearpage\part{The geometry of zeros}
\section{Growth and the geometry of the large zeros}\label{sec:zero_geometry}

\begin{readerroadmap}
The finite exponential representation gives a polygon of exponents. Large zeros can occur where two dominant terms nearly cancel. The proof first localizes those zeros and then uses symmetry and uniqueness to place them on the rays. Finite exceptional zeros remain part of the later products and sums.
\end{readerroadmap}
\modelfigure{zero_orbits}{The exact order-four model. The first three zeros on each of four rays are mapped by $w=z^4$ to three positive reduced zeros. The reduced axis uses a logarithmic scale. This complete order-four example does not classify the finite exceptional set at other orders.}{fig:zero-orbits}


We now assume $N\geq2$.  Write
\[
\eta_k=e^{(2k+1)\pi i/N},
\qquad 0\leq k<N.
\]
These are the vertices of a regular polygon and
\[
\Dfun_N(z)=\frac1N\sum_{k=0}^{N-1}e^{\eta_kz}.
\]

\begin{standardfact}[Order from Taylor coefficients]\label{sf:coefficient_order}
If $f(w)=\sum_{n\geq0}a_nw^n$ is an entire function with infinitely many nonzero coefficients, then
\[
\ord f
=
\limsup_{\substack{n\to\infty\\a_n\ne0}}
\frac{n\log n}{\log(1/|a_n|)}.
\]
Canonical products for entire functions of order below one have genus zero.  We use these standard facts in the form given in \cite{Boas1954,Titchmarsh1939}.
\end{standardfact}

\begin{proposition}[Orders and exponential type]\label{prop:orders}
For $N\geq2$,
\[
\ord\Delta_N=\frac1N.
\]
The function $\Dfun_N$ has order one and exponential type one.
\end{proposition}

\begin{proof}
The coefficient of $w^n$ in $\Delta_N$ has modulus $1/(Nn)!$.
Stirling's formula and \cref{sf:coefficient_order} give order $1/N$.
Since $M_{f(\,\cdot\,^N)}(r)=M_f(r^N)$ for the maximum modulus,
composition with $z^N$ multiplies the order by $N$ and gives order one
for $\Dfun_N$.

The finite exponential representation gives
\[
|\Dfun_N(z)|\leq e^{|z|},
\]
so the type is at most one.  Choose a root $\eta_0$ and put $z=t\overline{\eta_0}$ with $t>0$.  The term $e^{\eta_0z}=e^t$ has strictly larger real exponent than every other term.  Hence
\[
\Dfun_N(t\overline{\eta_0})
=
\frac1Ne^t(1+o(1)).
\]
The type is therefore one.
\end{proof}


\begin{definition}\label{def:indicator_polygon}
For $\theta\in\mathbb R$, define
\[
h_N(\theta)
=
\max_{0\leq k<N}
\cos\left(\theta+\frac{(2k+1)\pi}{N}\right).
\]
The tie directions are
\[
\mathfrak T_N
=
\left\{\frac{2\pi k}{N}:0\leq k<N\right\}
\quad\bmod 2\pi.
\]
\end{definition}

\begin{theorem}[Indicator polygon and dominant sectors]\label{thm:indicator_polygon}
The indicator of $\Dfun_N$ equals $h_N$, with the limsup taken over radii where $\Dfun_N(re^{i\theta})\ne0$:
\[
\limsup_{r\to\infty}
\frac1r\log|\Dfun_N(re^{i\theta})|
=
h_N(\theta).
\]
If $K$ is a compact set of directions disjoint from $\mathfrak T_N$, then there are $\delta>0$ and $R>0$ such that, for each $\theta\in K$, one root $\eta(\theta)$ of $-1$ satisfies
\[
\Dfun_N(re^{i\theta})
=
\frac1N
\exp\bigl(\eta(\theta)re^{i\theta}\bigr)
\left(1+O(e^{-\delta r})\right)
\]
uniformly for $r\geq R$.  In particular, no sufficiently large zero can lie in those directions.
\end{theorem}

\begin{proof}
For $z=re^{i\theta}$, the real part of the exponent $\eta_kz$ is
\[
r\cos\left(\theta+\frac{(2k+1)\pi}{N}\right).
\]
The triangle inequality gives the upper indicator bound $h_N(\theta)$.

Away from a tie direction, the maximizing vertex is unique.  It is constant
on each connected component of the complement of the tie directions.  A
compact set avoiding all ties meets only finitely many such components, and
compactness gives one uniform positive lower bound for every relevant
projection gap.  Division by the maximizing exponential gives the stated
uniform asymptotic and the zero-free conclusion.

At a tie direction exactly two adjacent vertices attain the maximum.  After rotation it is enough to take $\theta=0$.  With $c_N=\cos(\pi/N)$ and $s_N=\sin(\pi/N)$, the two dominant terms give
\[
\frac2N e^{c_Nr}\cos(s_Nr)
\]
and every other term is exponentially smaller.  Along a sequence on which $|\cos(s_Nr)|\geq1/2$, the lower indicator bound is $c_N=h_N(0)$.  Rotational symmetry treats every tie direction.
\end{proof}

The zero strings are governed by two adjacent exponents.  Put
\[
a_N=\frac\pi N,
\qquad
s_N=\sin a_N,
\qquad
c_N=\cos a_N,
\]
and
\[
A_N=\frac\pi{s_N},
\qquad
 t_m=A_N\left(m+\frac12\right).
\]
For $N\geq3$, also put
\[
\gamma_N=c_N-\cos(3a_N)>0.
\]

\begin{lemma}[Two term dominance]\label{lem:two_term_dominance}
Assume $N\geq3$.  There are positive constants $Y$, $R$, and $C$ with the following properties.

If $z=x+iy$ lies in the sector $|\arg z|\leq a_N$, has $x\geq R$, and is a zero of $\Dfun_N$, then $|y|\leq Y$.

On the strip $x\geq R$ and $|y|\leq Y$, define
\[
R_N(z)
=
\sum_{\substack{\eta^N=-1\\
\eta\ne e^{ia_N},e^{-ia_N}}}
\exp\bigl((\eta-e^{-ia_N})z\bigr).
\]
Then
\[
N e^{-e^{-ia_N}z}\Dfun_N(z)
=
1+e^{2is_Nz}+R_N(z),
\]
and, for $q=0,1,2$,
\[
|R_N^{(q)}(z)|
\leq
C e^{-\gamma_Nx}.
\]
\end{lemma}

\begin{proof}
A direct projection estimate gives the stronger bound
\[
Y=\frac{\log(N-1)}{2s_N}
\]
for every zero in the sector, without a large-radius assumption.
Write $z=re^{i\theta}$ with $0\leq\theta\leq a_N$ and $r\geq0$.
For $0\leq\ell\leq N-2$, the angle
$\theta+(2\ell+1)a_N$ lies between $a_N+\theta$ and
$2\pi-a_N-\theta$. Monotonicity of cosine on $[0,\pi]$, followed by
reflection about $\pi$, therefore gives
\[
\Re(\eta_\ell z)\leq r\cos(a_N+\theta).
\]
The remaining vertex $e^{-ia_N}$ has projection $r\cos(a_N-\theta)$.
Their difference is at least $2s_Nr\sin\theta=2s_Ny$.
If $\Dfun_N(z)=0$, divide its finite exponential sum by the term
with exponent $e^{-ia_N}z$ and apply the triangle inequality:
\[
1\leq(N-1)e^{-2s_Ny}.
\]
Thus $y\leq Y$. Conjugation gives the same bound for $-y$ on the lower
half sector. Since $N\geq3$ and $s_N>0$, this $Y$ is positive.

On the resulting strip, the roots $e^{ia_N}$ and $e^{-ia_N}$ have
real part $c_N$, while each remaining root $\eta=u+iv$ satisfies
$u\leq\cos(3a_N)=c_N-\gamma_N$. Hence
\[
\Re\bigl((\eta-e^{-ia_N})z\bigr)
\leq-\gamma_Nx+2Y.
\]
Division by $e^{e^{-ia_N}z}$ gives the displayed normalization.
Each differentiation multiplies a remainder term by
$\eta-e^{-ia_N}$, whose modulus is at most two. Thus, for $q=0,1,2$,
\[
|R_N^{(q)}(z)|\leq(N-2)2^qe^{2Y}e^{-\gamma_Nx}.
\]
One may take $R=1$ and $C=4(N-2)e^{2Y}$.
\end{proof}

\begin{lemma}[Near cancellation of two exponentials]\label{lem:near_cancellation}
For every $M>0$, there are constants $\varepsilon_M>0$ and $C_M>0$ such that the following holds.  If
\[
|\Re w|\leq M
\]
and
\[
|1+e^w|\leq\varepsilon_M,
\]
then there is a unique integer $m$ with
\[
|w-(2m+1)\pi i|
\leq
C_M|1+e^w|.
\]
\end{lemma}

\begin{proof}
In fact, the constants $\varepsilon_M=1/2$ and $C_M=3/2$ work
for every complex $w$, without the strip restriction. Put
$u=-(1+e^w)$ and $L=\log(1+u)=\log(-e^w)$, using the principal
logarithm. The elementary logarithm estimate on $|u|\leq1/2$ gives
$|L|\leq(3/2)|u|$. Since $e^{L+\pi i}=e^w$, there is an integer
$m$ such that
\[
w=L+(2m+1)\pi i.
\]
This gives the required distance bound. Distinct centers are separated
by at least $2\pi$, whereas two centers satisfying the bound would
be separated by at most $3|1+e^w|\leq3/2$. Thus the integer is unique.
The formal proof uses the convergent logarithm estimate and the exact
periodicity theorem for the complex exponential.
\end{proof}

\begin{theorem}[Large zero strings]\label{thm:large_zero_strings}
Assume $N\geq3$.  There are constants $m_0$, $C$, and $\gamma>0$ and a sequence of positive real simple zeros $(\tau_m)_{m\geq m_0}$ of $\Dfun_N$ such that
\[
\tau_m
=
 t_m+O(e^{-\gamma t_m}).
\]
Every zero of sufficiently large modulus is of the form
\[
\omega^k\tau_m
\]
for a unique $0\leq k<N$ and a unique $m\geq m_0$.

The derivatives satisfy
\[
\Dfun_N'(\omega^k\tau_m)
=
\omega^{-k}\Dfun_N'(\tau_m)
\]
and
\[
|\Dfun_N'(\tau_m)|
=
\frac{2s_N}{N}e^{c_N\tau_m}
\left(1+O(e^{-\gamma\tau_m})\right).
\]
\end{theorem}

\begin{proof}
Work first in the sector $|\arg z|\leq a_N$.  By \cref{lem:two_term_dominance}, every sufficiently large zero lies in a fixed strip about the positive real axis.  Put
\[
G(z)=N e^{-e^{-ia_N}z}\Dfun_N(z).
\]
Then
\[
G(z)=1+e^{2is_Nz}+R_N(z)
\]
with an exponentially small remainder and derivative.

If $G(z)=0$, then
\[
|1+e^{2is_Nz}|
\leq
C e^{-\gamma_N\Re z}.
\]
The real part of $2is_Nz$ is bounded in the strip.  Apply
\cref{lem:near_cancellation} with $w=2is_Nz$.  There is a unique
integer $m$ such that
\[
|z-t_m|\leq C_1e^{-\gamma_N\Re z},
\]
where
\[
 t_m=\frac{(2m+1)\pi}{2s_N}.
\]
In particular, $|\Re z-t_m|=o(1)$.  For all sufficiently large zeros
we have $|\Re z-t_m|\leq1$, and hence
\[
|z-t_m|\leq C_2e^{-\gamma_Nt_m}
\]
with a constant independent of $m$.
Because $\Re z\to+\infty$ in this argument, the corresponding integers
$m$ are eventually nonnegative and tend to infinity.

Fix a large $m$, put $a=2is_N$, and consider the closed disk
$B_m=\{z:|z-t_m|\leq 1/(2s_N)\}$. The remainder bounds hold
uniformly on these fixed-radius disks. For sufficiently large $m$,
$|R_N|\leq 1/2$ and $|R_N'|\leq s_N/2$ there. Define
\[
\Phi_m(z)=t_m+\frac{\log(1+R_N(z))}{2is_N},
\]
using the principal logarithm. The estimate
$|\log(1+u)|\leq(3/2)|u|$ shows that $\Phi_m$ maps $B_m$ into
itself. The derivative bound $|\Phi_m'|\leq1/2$ makes it a contraction.
Banach's fixed-point theorem gives a unique fixed point.
On $B_m$, $|a(z-t_m)|\leq1<\pi$, so the principal logarithm inverts
the exponential and the fixed-point equation is equivalent to $G(z)=0$.
Moreover,
\[
|z-t_m|\leq\frac{3}{4s_N}|R_N(z)|=O(e^{-\gamma_Nt_m}).
\]
At this zero,
\[
|G'(z)|\geq |a|\,|1+R_N(z)|-|R_N'(z)|\geq s_N/2>0,
\]
so the zero is simple. The disks are mutually disjoint because their
diameter $1/s_N$ is smaller than the center spacing $\pi/s_N$.
Enlarging the strip if necessary, these disks lie inside the fundamental
sector and the fixed strip for large $m$.

The localization estimate puts every sufficiently large zero in the
fundamental sector into one of these disks.  The zero sets of $G$ and
$\Dfun_N$ coincide because their prefactor is nonzero.  Each disk is
invariant under complex conjugation, and $\Dfun_N$ has real
coefficients.  The conjugate of its unique zero is therefore the same
zero.  It is real, positive for large $m$, and simple by the
derivative bound above.  Denote it by $\tau_m$.

The parts of the fundamental sector away from the positive tie ray,
including its boundary rays, were already excluded uniformly.
Rotational invariance now gives every remaining large zero as
$\omega^k\tau_m$.  Distinct rays and the disjoint disks give uniqueness
of $k$ and $m$.  Taking $\gamma=\gamma_N$ proves the zero location in
the statement.

Differentiate the normalized expression.  At $z=\tau_m$,
\[
G'(\tau_m)
=
2is_Ne^{2is_N\tau_m}+R_N'(\tau_m)
=
-2is_N+O(e^{-\gamma_N\tau_m}).
\]
Since $G(\tau_m)=0$,
\[
\Dfun_N'(\tau_m)
=
\frac1N e^{e^{-ia_N}\tau_m}G'(\tau_m).
\]
Taking absolute values gives the derivative asymptotic.  Differentiating $\Dfun_N(\omega z)=\Dfun_N(z)$ gives the rotation formula.
\end{proof}

\begin{theorem}[Leading displacement term and derivative sign]\label{thm:first_zero_correction}
With the notation of \cref{lem:two_term_dominance,thm:large_zero_strings},
\[
\tau_m
=
t_m+
\frac{R_N(t_m)}{2is_N}
+
O(e^{-2\gamma_Nt_m}).
\]
The quotient $R_N(t_m)/(2is_N)$ is real.  One also has
\[
\Dfun_N'(\tau_m)
=
(-1)^{m+1}
\frac{2s_N}{N}
e^{c_N\tau_m}
\left(1+O(e^{-\gamma_N\tau_m})\right).
\]
\end{theorem}

\begin{proof}
Put $\delta_m=\tau_m-t_m$.  The large zero theorem gives $\delta_m=O(e^{-\gamma_Nt_m})$.  Since $e^{2is_Nt_m}=-1$,
\[
1+e^{2is_N(t_m+\delta_m)}
=
-2is_N\delta_m+O(\delta_m^2).
\]
The remainder estimate through the first derivative gives
\[
R_N(t_m+\delta_m)
=
R_N(t_m)+O(e^{-\gamma_Nt_m}|\delta_m|).
\]
Substitution in the zero equation gives the displacement formula.

At the real point $t_m$, the normalized value
\[
N e^{-e^{-ia_N}t_m}\Dfun_N(t_m)
\]
is purely imaginary.  The two dominant terms cancel there, so $R_N(t_m)$ is purely imaginary and the displayed correction is real.

For the derivative, use the formula in the proof of \cref{thm:large_zero_strings} and
\[
e^{-is_N\tau_m}
=
(-1)^m(-i)
\left(1+O(e^{-\gamma_N\tau_m})\right).
\]
The stated real sign follows.
\end{proof}

\begin{corollary}[Leading correction for $N=3$]\label{cor:N3_zero_correction}
For $N=3$,
\[
t_m=\frac{2\pi}{\sqrt3}\left(m+\frac12\right)
\]
and
\[
\tau_m
=
t_m+
\frac{(-1)^m}{\sqrt3}e^{-3t_m/2}
+
O(e^{-3t_m}).
\]
\end{corollary}

\begin{proof}
The only nondominant exponent is $-1$.  Hence
\[
R_3(t)
=
\exp\left(\left(-\frac32+\frac{\sqrt3}{2}i\right)t\right).
\]
At $t=t_m$, its phase is $i(-1)^m$.  Insert this value into \cref{thm:first_zero_correction}.
\end{proof}

\begin{corollary}[The case $N=2$]\label{cor:case_two_zeros}
For $N=2$,
\[
\Dfun_2(z)=\cos z.
\]
Its zeros are
\[
\pm\left(m+\frac12\right)\pi,
\qquad m\geq0,
\]
and all are simple.  They agree with the formula $t_m=A_2(m+1/2)$.
\end{corollary}

\begin{proof}
This is the standard zero set of the cosine function.
\end{proof}

\begin{corollary}[Eventual simplicity and exact rays]\label{cor:eventual_simplicity}
For every $N\geq2$, all but finitely many zeros of $\Dfun_N$ are simple.  Every sufficiently large zero lies on one of the rays
\[
\omega^k(0,\infty),
\qquad 0\leq k<N.
\]
\end{corollary}

\begin{proof}
Use \cref{thm:large_zero_strings} for $N\geq3$ and \cref{cor:case_two_zeros} for $N=2$.
\end{proof}


\begin{example}[Exact fourth-order strings]\label{ex:fourth_order_zeros}
For $N=4$,
\[
\Dfun_4(z)
=
\cos\left(\frac z{\sqrt2}\right)
\cosh\left(\frac z{\sqrt2}\right).
\]
The positive zeros are exactly
\[
\tau_m
=
\sqrt2\pi\left(m+\frac12\right)
=
A_4\left(m+\frac12\right).
\]
Rotation by the fourth roots of unity gives the real and imaginary zero strings.  In this case the error term in \cref{thm:large_zero_strings} is identically zero.
\end{example}

\begin{corollary}[Reduced zeros]\label{cor:reduced_zeros}
Let $\Lambda_N$ be the nonzero zero multiset of $\Delta_N$.  All but finitely many members of $\Lambda_N$ are positive and simple.  They can be indexed as
\[
\lambda_m=\tau_m^N
\]
for all sufficiently large $m$.  If $N\geq3$, the notation may be chosen
as in \cref{thm:large_zero_strings}, and
\[
\lambda_m
=
 t_m^N
+O\left(t_m^{N-1}e^{-\gamma t_m}\right).
\]
For $N=2$, take $\tau_m=t_m=\pi(m+1/2)$ for $m\geq0$; then
$\lambda_m=t_m^2$ exactly.
\end{corollary}

\begin{proof}
The identity $\Dfun_N(z)=\Delta_N(z^N)$ identifies each orbit of $N$ nonzero zeros of $\Dfun_N$ with one zero of $\Delta_N$.  For $N\geq3$, the large orbit representatives are the positive zeros $\tau_m$ from \cref{thm:large_zero_strings}; for $N=2$, use \cref{cor:case_two_zeros}.  Simplicity is preserved because
\[
\Dfun_N'(z)=Nz^{N-1}\Delta_N'(z^N)
\]
and $z\ne0$.  The asymptotic follows by expanding $\tau_m^N$ about $t_m$.
\end{proof}

\begin{theorem}[Zero counting asymptotics]\label{thm:zero_counting}
Let $n_{\Dfun_N}(R)$ count the zeros of $\Dfun_N$ in $|z|\leq R$, with multiplicity.  Let $n_{\Delta_N}(R)$ count the zeros of $\Delta_N$ in $|w|\leq R$.  Then, as $R\to\infty$,
\[
n_{\Dfun_N}(R)
=
\frac{N\sin(\pi/N)}{\pi}R+O(1)
\]
and
\[
n_{\Delta_N}(R)
=
\frac{\sin(\pi/N)}{\pi}R^{1/N}+O(1).
\]
\end{theorem}

\begin{proof}
Each of the $N$ rays contains one large zero near
\[
 t_m=\frac\pi{\sin(\pi/N)}\left(m+\frac12\right).
\]
Since $\tau_m-t_m=o(1)$ and the model spacing is the positive constant
$A_N$, replacing $\tau_m\leq R$ by $t_m\leq R$ changes the count by at
most a bounded amount.
The number of such indices with $t_m\leq R$ is
\[
\frac{\sin(\pi/N)}\pi R+O(1).
\]
Multiplication by $N$ and addition of the finite exceptional set give the first formula.  The second follows from
\[
n_{\Dfun_N}(R^{1/N})=Nn_{\Delta_N}(R).
\]
\end{proof}


\begin{corollary}[Zero exponents]\label{cor:exponent_of_convergence}
The exponent of convergence of the nonzero zeros of $\Dfun_N$ is one.  The exponent of convergence of the zeros of $\Delta_N$ is $1/N$.
\end{corollary}

\begin{proof}
The asymptotic counting formulas in \cref{thm:zero_counting} give the two exponents by partial summation.
\end{proof}

\clearpage\part{Products, moments, and zero zeta functions}
\section{Hadamard products and zero power sums}\label{sec:hadamard}

\begin{keyidea}
The product uses every reduced zero with its multiplicity. Its logarithm turns multiplication into addition, so reciprocal-zero power sums become coefficients. This step depends on the global product theorem. A truncated zero list provides a numerical partial sum, while the identity uses the entire multiset.
\end{keyidea}
\modelfigure{zero_moments}{Order-two moment sums over the first forty reduced zeros $\lambda_k=\pi^2(k+1/2)^2$. The curves are finite numerical partial sums. The dashed limits $1/2$ and $1/6$ follow from the exact logarithmic coefficients and the product identity.}{fig:zero-moments}


The reduction from $z$ to $w=z^N$ removes the rotational repetition of the zeros.

\begin{theorem}[Genus zero product]\label{thm:genus_zero_product}
For $N\geq2$, let $\Lambda_N$ be the nonzero zero multiset of $\Delta_N$.  Then
\[
\Delta_N(w)
=
\prod_{\lambda\in\Lambda_N}
\left(1-\frac w\lambda\right),
\]
where the product converges locally uniformly.  Consequently,
\[
\Dfun_N(z)
=
\prod_{\lambda\in\Lambda_N}
\left(1-\frac{z^N}{\lambda}\right).
\]
\end{theorem}

\begin{proof}
By \cref{prop:orders}, the order of $\Delta_N$ is $1/N<1$.  By
\cref{cor:reduced_zeros}, it has infinitely many zeros, and
\cref{cor:exponent_of_convergence} gives
\[
\sum_{\lambda\in\Lambda_N}|\lambda|^{-1}<\infty.
\]
Hadamard factorization therefore gives a canonical product of genus zero.
The only possible exponential factor is a nonzero constant because the order
is below one.  Since $\Delta_N(0)=1$, that constant is one.  Substitution of
$w=z^N$ gives the second formula.
\end{proof}


\begin{theorem}[Diagonal Fredholm determinant representation]\label{thm:trace_class_model}
Let $I_N$ be a multiplicity-sensitive countable index set for $\Lambda_N$
and put $\mathcal H_N=\ell^2(I_N)$.  On its standard orthonormal basis,
define the diagonal operator
\[
K_Ne_\lambda=\lambda^{-1}e_\lambda.
\]
Then $K_N$ is trace class and
\[
\Delta_N(w)=\Fdet(I-wK_N).
\]
For every $m\geq1$,
\[
\kappa_{N,m}=\tr(K_N^m).
\]
\end{theorem}

\begin{proof}
By \cref{cor:exponent_of_convergence},
\[
\sum_{\lambda\in\Lambda_N}|\lambda|^{-1}<\infty.
\]
Thus the diagonal operator $K_N$ is trace class.  Its Fredholm determinant is the locally uniform product
\[
\Fdet(I-wK_N)
=
\prod_{\lambda\in\Lambda_N}
\left(1-\frac w\lambda\right),
\]
which equals $\Delta_N(w)$ by \cref{thm:genus_zero_product}.

Taking the logarithm gives
\[
-\log\Delta_N(w)
=
\sum_{m\geq1}\frac{w^m}{m}\tr(K_N^m)
\]
near the origin.  Comparison with the definition of $\kappa_{N,m}$ proves the trace identity.
\end{proof}

This is the standard diagonal encoding of this family's genus zero product.
It is not a general trace-class resolvent or Feshbach theorem.

\begin{corollary}[Complete symmetric function form]\label{cor:complete_symmetric_form}
For a positive integer $r$, let each eigenvalue $\lambda^{-1}$ of $K_N$ be repeated $r$ times.  If $h_n$ denotes the complete homogeneous symmetric function of degree $n$ in this absolutely summable multiset, then
\[
\frac{P_{N,n}(r)}{(Nn)!}=h_n.
\]
For arbitrary complex $r$, the Bell polynomial formula is the polynomial continuation obtained by multiplying every power sum by $r$.
\end{corollary}

\begin{proof}
The generating function for the complete homogeneous symmetric functions is
\[
\prod_{\lambda\in\Lambda_N}
\left(1-\frac w\lambda\right)^{-r}
=
\Delta_N(w)^{-r}.
\]
Coefficient comparison gives the first identity.  Power sums determine the complete homogeneous functions through the Newton recurrence.  Replacing each power sum by $r$ times itself gives \cref{thm:reciprocal_bell,thm:logarithmic_newton_recurrence}.
\end{proof}


\begin{proposition}[Products for the normalized sections]\label{prop:normalized_section_product}
For every $0\leq j<N$, the entire function
\[
\Phi_{N,j}(w)
=
j!\sum_{n\geq0}\frac{w^n}{(Nn+j)!}
\]
has order $1/N$.  If $N\geq2$ and $\Lambda_{N,j}$ is its nonzero zero multiset, then
\[
\Phi_{N,j}(w)
=
\prod_{\mu\in\Lambda_{N,j}}
\left(1-\frac w\mu\right).
\]
For $m\geq1$, put
\[
\sigma_{N,j,m}
=
\sum_{\mu\in\Lambda_{N,j}}\mu^{-m}.
\]
Then
\[
C_{N,j,n}(r)
=
\frac{(Nn)!}{n!}
B_n\bigl(
 r\sigma_{N,j,1},
 r1!\sigma_{N,j,2},
 \ldots,
 r(n-1)!\sigma_{N,j,n}
\bigr).
\]
\end{proposition}

\begin{proof}
The coefficient formula for the order gives $1/N$.  The function is not a polynomial.  An entire function of order below one with only finitely many zeros would be a polynomial times a constant by Hadamard factorization, which is impossible here.  The genus zero product follows as in \cref{thm:genus_zero_product}.

Taking the logarithm gives
\[
-\log\Phi_{N,j}(w)
=
\sum_{m\geq1}\sigma_{N,j,m}\frac{w^m}{m}.
\]
Since $\Cfun_{N,j}(z)^r=\Phi_{N,j}(z^N)^{-r}$, the Bell polynomial calculation from \cref{thm:reciprocal_bell} gives the coefficient formula.
\end{proof}

\begin{definition}\label{def:zero_power_sums}
For $m\geq1$, define
\[
\sigma_{N,m}
=
\sum_{\lambda\in\Lambda_N}\lambda^{-m},
\]
with multiplicities.
\end{definition}

The series converges absolutely because the exponent of convergence of $\Lambda_N$ is $1/N<1$.

\begin{theorem}[Newton sums for the reduced zeros]\label{thm:newton_sums}
For $|w|$ below the least zero modulus, with the analytic logarithm
normalized to zero at the origin,
\[
-\log\Delta_N(w)
=
\sum_{m\geq1}\sigma_{N,m}\frac{w^m}{m}.
\]
Thus
\[
\kappa_{N,m}=\sigma_{N,m}.
\]
If
\[
a_m=\frac{(-1)^m}{(Nm)!},
\qquad a_0=1,
\]
then
\[
\sigma_{N,m}
=
-ma_m-
\sum_{k=1}^{m-1}a_k\sigma_{N,m-k}.
\]
\end{theorem}

\begin{proof}
Take the logarithm of the product in \cref{thm:genus_zero_product}.  Absolute convergence in a disk below the first zero permits termwise expansion
\[
-\log\left(1-\frac w\lambda\right)
=
\sum_{m\geq1}\frac{w^m}{m\lambda^m}.
\]
This gives the logarithmic formula and identifies $\kappa_{N,m}$ with $\sigma_{N,m}$.

Differentiate the product identity in logarithmic form to obtain
\[
\Delta_N'(w)
=
-\Delta_N(w)
\sum_{m\geq1}\sigma_{N,m}w^{m-1}.
\]
Comparison of the coefficient of $w^{m-1}$ gives
\[
ma_m
=
-\sum_{k=0}^{m-1}a_k\sigma_{N,m-k}.
\]
Solving for $\sigma_{N,m}$ gives the recurrence.
\end{proof}



\begin{corollary}[Convergent Mittag--Leffler partial-fraction form]\label{cor:mittag_leffler_tangent}
Away from the zeros of $\Dfun_N$,
\[
\Tfun_N(z)
=
Nz^{N-1}
\sum_{\lambda\in\Lambda_N}
\frac1{\lambda-z^N}.
\]
The series converges absolutely and locally uniformly on compact sets
away from the poles.  Thus no regularizing terms are required.
\end{corollary}

\begin{proof}
Logarithmic differentiation of the genus zero product gives
\[
-\frac{\Delta_N'(w)}{\Delta_N(w)}
=
\sum_{\lambda\in\Lambda_N}\frac1{\lambda-w}.
\]
The series is locally uniform because $\sum_{\lambda}|\lambda|^{-1}<\infty$.  Substitute $w=z^N$ and use the chain rule.
\end{proof}

\begin{proposition}[Theta Bell polynomial formula]\label{prop:theta_bell}
Put
\[
\widetilde{\sigma}_{N,m}=(-1)^m\sigma_{N,m}.
\]
Then
\[
\Thetafun_{N,n}(u)
=
\frac{(Nn)!}{n!}
B_n\bigl(
 \widetilde{\sigma}_{N,1}(u^N+u),
 1!\widetilde{\sigma}_{N,2}(u^{2N}+u),
 \ldots,
 (n-1)!\widetilde{\sigma}_{N,n}(u^{nN}+u)
\bigr).
\]
\end{proposition}

\begin{proof}
The reduced unsigned denominator is
\[
\Psi_N(w)=\sum_{n\geq0}\frac{w^n}{(Nn)!}=\Delta_N(-w).
\]
By \cref{thm:newton_sums},
\[
-\log\Psi_N(w)
=
\sum_{m\geq1}\widetilde{\sigma}_{N,m}\frac{w^m}{m}.
\]
With $w=z^N$, the logarithm of the theta generating function is
\[
\sum_{m\geq1}
\widetilde{\sigma}_{N,m}(u^{Nm}+u)\frac{w^m}{m}.
\]
The complete Bell polynomial formula follows.
\end{proof}

\begin{corollary}[Tangent coefficients as zero sums]\label{cor:tangent_zero_sums}
For every $n\geq0$,
\[
T_N(n)
=
N(Nn+N-1)!\sigma_{N,n+1}.
\]
\end{corollary}

\begin{proof}
Combine \cref{prop:tangent_kappa,thm:newton_sums}.
\end{proof}

\begin{corollary}[Reciprocal coefficients from zero data]\label{cor:reciprocal_zero_data}
For every $n\geq0$,
\[
P_{N,n}(r)
=
\frac{(Nn)!}{n!}
B_n\bigl(
 r\sigma_{N,1},
 r1!\sigma_{N,2},
 \ldots,
 r(n-1)!\sigma_{N,n}
\bigr).
\]
\end{corollary}

\begin{proof}
Use \cref{thm:reciprocal_bell,thm:newton_sums}.
\end{proof}

\begin{example}[The cosine orbit sums]\label{ex:cosine_orbit_sums}
For $N=2$,
\[
\Delta_2(w)=\cos\sqrt w.
\]
Its zeros are
\[
\lambda_m=\pi^2\left(m+\frac12\right)^2,
\qquad m\geq0.
\]
Hence
\[
\sigma_{2,q}
=
\frac{2^{2q}-1}{\pi^{2q}}\zeta(2q).
\]
The identity in \cref{cor:tangent_zero_sums} recovers the classical tangent coefficients from these zero power sums.
\end{example}

\section{Coefficient asymptotics from the nearest zeros}\label{sec:coefficient_asymptotics}

For an entire function $f$, let $\Zcal(f)$ denote its multiset of nonzero zeros.  Let $r$ be a positive integer.  Then $\Dfun_N^{-r}$ is meromorphic and a simple zero of $\Dfun_N$ becomes a pole of order $r$.

\begin{theorem}[Nearest-pole coefficient expansion]\label{thm:nearest_circle_expansion}
Let $r\in\ZZ_{>0}$.  Suppose the zeros of $\Dfun_N$ of least modulus
$R$ are simple and form the finite set $\Zcal_R$.  Suppose the next
zero modulus is $R_2>R$.  At each $\rho\in\Zcal_R$, write the principal
part
\[
\Dfun_N(z)^{-r}
=
\sum_{\ell=1}^{r}
 a_{\rho,\ell}
\left(1-\frac z\rho\right)^{-\ell}
+
O(1).
\]
Then, for every $R_*$ with $R<R_*<R_2$ and $m=Nn$, as $n\to\infty$,
\[
P_{N,n}(r)
=
 m!
\sum_{\rho\in\Zcal_R}
\sum_{\ell=1}^{r}
 a_{\rho,\ell}
\binom{m+\ell-1}{\ell-1}
\rho^{-m}
+
O_{R_*}(m!R_*^{-m}).
\]
The leading Laurent coefficient is
\[
a_{\rho,r}
=
(-\rho\Dfun_N'(\rho))^{-r}.
\]
\end{theorem}

\begin{proof}
Subtract the complete principal parts at the zeros on $|z|=R$.  The remainder is holomorphic on $|z|<R_2$, hence on a neighborhood of every closed disk $|z|\leq R_*<R_2$.  Cauchy's estimate gives an ordinary coefficient bound $O_{R_*}(R_*^{-m})$.  The coefficient of $z^m$ in
\[
\left(1-\frac z\rho\right)^{-\ell}
\]
is
\[
\binom{m+\ell-1}{\ell-1}\rho^{-m}.
\]
Multiplication by $m!$ converts the ordinary coefficient to $P_{N,n}(r)$.

Near a simple zero,
\[
\Dfun_N(z)
=
-\rho\Dfun_N'(\rho)
\left(1-\frac z\rho\right)
+O\left(\left(1-\frac z\rho\right)^2\right).
\]
Raising to the power $-r$ gives the formula for $a_{\rho,r}$.
\end{proof}

\begin{corollary}[Leading dominant zero term]\label{cor:leading_dominant_zero}
Under the hypotheses of \cref{thm:nearest_circle_expansion},
\[
\begin{aligned}
P_{N,n}(r)
={}&m!\binom{m+r-1}{r-1}
\sum_{\rho\in\Zcal_R}
\frac{\rho^{-m}}{(-\rho\Dfun_N'(\rho))^r}
\\
&+O\left(m!m^{r-2}R^{-m}\right)
+O_{R_*}\left(m!R_*^{-m}\right).
\end{aligned}
\]
For $r=1$, the first error term is absent.
\end{corollary}

\begin{proof}
The terms with $\ell<r$ in \cref{thm:nearest_circle_expansion} are bounded by a constant times $m^{r-2}R^{-m}$.  The term with $\ell=r$ gives the displayed main term.
\end{proof}


\begin{corollary}[Orbit reduced nearest circle term]\label{cor:orbit_reduced_asymptotic}
Assume the hypotheses of \cref{thm:nearest_circle_expansion}.  Choose one representative from each orbit of the nearest zeros under multiplication by $\omega$, and call the resulting set $\Zcal_R^{\mathrm{orb}}$.  With $m=Nn$,
\[
\begin{aligned}
P_{N,n}(r)
={}&N m!\binom{m+r-1}{r-1}
\sum_{\tau\in\Zcal_R^{\mathrm{orb}}}
\frac{\tau^{-m}}{(-\tau\Dfun_N'(\tau))^r}
\\
&+O\left(m!m^{r-2}R^{-m}\right)
+O_{R_*}\left(m!R_*^{-m}\right).
\end{aligned}
\]
For $r=1$, the first error term is absent.
\end{corollary}

\begin{proof}
If $\rho=\omega^k\tau$, then
\[
\Dfun_N'(\rho)=\omega^{-k}\Dfun_N'(\tau).
\]
Therefore
\[
-\rho\Dfun_N'(\rho)
=-\tau\Dfun_N'(\tau).
\]
Since $m=Nn$, one also has $\rho^{-m}=\tau^{-m}$.  Every orbit therefore contributes $N$ equal terms in \cref{cor:leading_dominant_zero}.
\end{proof}

\begin{corollary}[Secant asymptotic]\label{cor:secant_asymptotic}
For $N=2$ and $r=1$,
\[
P_{2,n}(1)
\sim
2(2n)!\left(\frac2\pi\right)^{2n+1}.
\]
\end{corollary}

\begin{proof}
The nearest zeros of $\cos z$ are $\rho=\pm\pi/2$, both simple.  Their contributions in \cref{cor:leading_dominant_zero} are equal.
\end{proof}

\section{Zero zeta functions}\label{sec:zero_zeta}

The assumption $N\geq2$ remains in force throughout this section.  The zero counting theorem gives the exact abscissae of absolute convergence.  The exponentially accurate zero locations give more.  They compare the zero zeta with a Hurwitz zeta function.

\begin{definition}\label{def:absolute_zero_zetas}
Define
\[
\mathfrak Z_N^{\mathrm{abs}}(s)
=
\sum_{\rho\in\Zcal(\Dfun_N)}|\rho|^{-s}
\]
and
\[
\mathfrak Z_N^{\mathrm{red}}(s)
=
\sum_{\lambda\in\Lambda_N}|\lambda|^{-s},
\]
with multiplicities.  Each term means $e^{-s\log|\rho|}$ or $e^{-s\log|\lambda|}$.
\end{definition}

\begin{theorem}[Abscissae of absolute convergence]\label{thm:zeta_convergence}
The series $\mathfrak Z_N^{\mathrm{abs}}$ converges absolutely and locally uniformly when
\[
\Re s>1.
\]
It fails to converge absolutely when $\Re s\leq1$.  The series $\mathfrak Z_N^{\mathrm{red}}$ converges absolutely and locally uniformly when
\[
\Re s>\frac1N.
\]
It fails to converge absolutely when $\Re s\leq1/N$.  Both functions are holomorphic on their open half planes of absolute convergence.
\end{theorem}

\begin{proof}
By \cref{thm:zero_counting}, the zeros of $\Dfun_N$ have linear counting density, while the reduced zeros have counting order $R^{1/N}$.  Partial summation gives absolute convergence in the stated open half planes.  The positive terms obtained after taking absolute values diverge on the boundary and to its left.  Local uniform convergence follows by the Weierstrass test on compact subsets.
\end{proof}

\begin{theorem}[Meromorphic continuation of the symmetry-orbit zero zeta]\label{thm:zeta_continuation}
Put
\[
A_N=\frac\pi{\sin(\pi/N)}.
\]
There is an entire function $Q_N^{\mathrm{red}}$ such that
\[
\mathfrak Z_N^{\mathrm{red}}(s)
=
A_N^{-Ns}\zeta\left(Ns,\frac12\right)+Q_N^{\mathrm{red}}(s)
\]
on $\Re s>1/N$.  The right hand side gives a meromorphic continuation to $\CC$ with one simple pole at $s=1/N$.  Its residue is
\[
\frac{\sin(\pi/N)}{N\pi}.
\]

There is an entire function $Q_N^{\mathrm{abs}}$ such that, initially on
$\Re s>1$,
\[
\mathfrak Z_N^{\mathrm{abs}}(s)
=
N A_N^{-s}\zeta\left(s,\frac12\right)+Q_N^{\mathrm{abs}}(s).
\]
Thus $\mathfrak Z_N^{\mathrm{abs}}$ continues meromorphically to $\CC$ with one simple pole at $s=1$, of residue
\[
\frac{N\sin(\pi/N)}\pi.
\]
\end{theorem}

\begin{proof}
For $N=2$, \cref{cor:case_two_zeros} and
$\Dfun_2(z)=\Delta_2(z^2)$ give
\[
\lambda_m=\pi^2\left(m+\frac12\right)^2,
\qquad m\geq0.
\]
Thus the first formula holds with $Q_2^{\mathrm{red}}=0$.

For $N\geq3$, \cref{cor:reduced_zeros} writes all but finitely many reduced zeros as
\[
\lambda_m
=
\left(A_N\left(m+\frac12\right)\right)^N(1+\varepsilon_m)
\]
where $\lambda_m>0$, $1+\varepsilon_m>0$, and
\[
\varepsilon_m
=
O\left(t_m^{-1}e^{-\gamma t_m}\right).
\]
On every compact set of $s$ values,
\[
(1+\varepsilon_m)^{-s}-1
=
O(|\varepsilon_m|)
\]
uniformly for large $m$.  The difference
\[
\lambda_m^{-s}
-
A_N^{-Ns}\left(m+\frac12\right)^{-Ns}
\]
therefore has a normally convergent sum on every compact subset of
$\CC$.  Indeed, the model power has at most polynomial growth in $m$
on a compact set, while $|\varepsilon_m|$ decays exponentially.  The
finitely many exceptional terms $|\lambda|^{-s}$ and the finitely many
omitted model terms are entire functions of $s$.  Their addition to the
normally convergent tail defines $Q_N^{\mathrm{red}}$ and proves the
first formula. Thus that formula is established for every $N\geq2$.

In either case, the Hurwitz zeta function has a meromorphic continuation with one simple pole at its first argument equal to one and residue one.  Composition with $Ns$ gives the stated pole and residue.

For every $N\geq2$, each nonzero zero $\lambda$ of $\Delta_N$ gives $N$ zeros of $\Dfun_N$, each of modulus $|\lambda|^{1/N}$.  Hence
\[
\mathfrak Z_N^{\mathrm{abs}}(s)
=
N\mathfrak Z_N^{\mathrm{red}}(s/N).
\]
Substitution gives the second formula with
\[
Q_N^{\mathrm{abs}}(s)=NQ_N^{\mathrm{red}}(s/N),
\]
which is entire.
\end{proof}

\begin{corollary}[Exact zeta models in orders two and four]\label{cor:exact_zeta_models}
For $N=2$,
\[
\mathfrak Z_2^{\mathrm{red}}(s)
=
\pi^{-2s}\zeta\left(2s,\frac12\right)
\]
and
\[
\mathfrak Z_2^{\mathrm{abs}}(s)
=
2\pi^{-s}\zeta\left(s,\frac12\right).
\]
For $N=4$, put $A_4=\sqrt2\pi$.  Then
\[
\mathfrak Z_4^{\mathrm{red}}(s)
=
A_4^{-4s}\zeta\left(4s,\frac12\right)
\]
and
\[
\mathfrak Z_4^{\mathrm{abs}}(s)
=
4A_4^{-s}\zeta\left(s,\frac12\right).
\]
In these two cases the entire correction terms in \cref{thm:zeta_continuation} vanish identically.
\end{corollary}

\begin{proof}
For $N=2$, the reduced zeros are
\[
\pi^2\left(m+\frac12\right)^2,
\qquad m\geq0.
\]
For $N=4$, \cref{ex:fourth_order_zeros} gives the reduced zeros
\[
A_4^4\left(m+\frac12\right)^4,
\qquad m\geq0.
\]
Summing their negative powers proves the reduced formulas.  The absolute formulas follow from
\[
\mathfrak Z_N^{\mathrm{abs}}(s)
=
N\mathfrak Z_N^{\mathrm{red}}(s/N).
\]
\end{proof}

\begin{definition}\label{def:branch_orbit_zeta}
Choose a logarithm for each member of the finite exceptional part of $\Lambda_N$.  Use the real logarithm for the eventual positive sequence.  Define
\[
\mathfrak Z_N^{\mathrm{orb}}(s)
=
\sum_{\lambda\in\Lambda_N}e^{-s\Log\lambda}.
\]
This series is initially defined on $\Re s>1/N$.
\end{definition}

\begin{corollary}[Branch-dependent symmetry-orbit continuation]\label{cor:branch_grouped_continuation}
The branch-dependent orbit zeta $\mathfrak Z_N^{\mathrm{orb}}$ continues meromorphically to $\CC$.  It has the same single pole and residue as $\mathfrak Z_N^{\mathrm{red}}$.  Changing the logarithms of the finite exceptional zeros changes the continuation by an entire function.
\end{corollary}

\begin{proof}
The tail consists of positive reduced zeros, so it has the same comparison with the Hurwitz zeta.  Only finitely many terms depend on the branch choices.
\end{proof}

\begin{definition}\label{def:derivative_weighted_zeta}
Let $q$ be a positive integer.  Choose logarithms of the simple zeros of $\Dfun_N$ with arguments in a fixed bounded interval.  Define
\[
\mathfrak Z_{N,q}^{\mathrm{der}}(s)
=
\sum_{\substack{\Dfun_N(\rho)=0\\ \Dfun_N'(\rho)\ne0}}
\frac{e^{-s\Log\rho}}{\Dfun_N'(\rho)^q}.
\]
\end{definition}

\begin{theorem}[Derivative-weighted convergence]\label{thm:derivative_weighted_convergence}
If $N\geq3$ and $q\geq1$, the series $\mathfrak Z_{N,q}^{\mathrm{der}}(s)$ converges locally uniformly for every $s\in\CC$.  It therefore defines an entire function.

For $N=2$, the series converges absolutely and locally uniformly on
$\Re s>1$, where it defines a holomorphic function.
\end{theorem}

\begin{proof}
For $N\geq3$, \cref{thm:large_zero_strings} gives
\[
|\Dfun_N'(\rho)|^{-1}
\leq
C e^{-c_N|\rho|}
\]
for every sufficiently large simple zero.  Let $L$ be a compact subset
of the $s$ plane.  Because the chosen arguments of the zeros lie in a
fixed bounded interval, there are constants $C_L$ and $M_L$ such that
\[
\sup_{s\in L}|e^{-s\Log\rho}|
\leq C_L|\rho|^{M_L}
\]
for every sufficiently large zero.  Hence the absolute value of the
corresponding summand is at most
\[
C_L'|\rho|^{M_L}e^{-q c_N|\rho|}.
\]
The linear zero counting estimate makes this majorant summable.  The
Weierstrass test gives absolute local uniform convergence of the
series of entire summands, so its sum is entire.  The finite exceptional
set causes no difficulty.

For $N=2$, the derivative of the cosine has modulus one at every zero.  The assertion then follows from the linear zero spacing.
\end{proof}

\begin{remark}\label{rem:multiple_zeros}
The derivative weighted sum omits the finitely many zeros that might be multiple.  A treatment of those zeros requires the complete principal parts of $\Dfun_N^{-q}$ rather than the single derivative weight.
\end{remark}

\clearpage\part{Scope and formal evidence}
\section{Scope boundaries and further questions}\label{sec:further_questions}

The large-zero theorem settles eventual simplicity but does not classify the
finite exceptional set. A remaining question is whether every zero of
$\Dfun_N$ is simple and every reduced zero of $\Delta_N$ is positive. Both
statements hold for $N=2$ and $N=4$. The argument here proves them only outside
a compact set for general $N$.

The leading correction in \cref{thm:first_zero_correction} may admit iteration
to a full exponential expansion for $\tau_m$. Such an expansion is an open
problem, not a result of this paper. It could refine the entire correction in
the zero-zeta comparison and lead to explicit special values.

The symmetry-orbit zero zeta has a meromorphic continuation with one pole. No
functional equation is obtained. A branch dependent sum over the $N$ roots of
each reduced zero must specify logarithms on the rays.

The coefficient families in
\cref{sec:reciprocal_powers,sec:residue_section,sec:theta} have exact formal
parameter-addition laws. Their arithmetic behavior, including odd-prime
congruences, dyadic weighting, and $p$-adic or Amice--Mahler interpolation,
belongs to Paper IX. General bounded-operator semigroup, resolvent, and
Feshbach theory belongs to Paper VI. Paper III owns general nonlinear
operation trees, while abstract completion and occupancy questions belong to
Paper X. None of those theories is inferred from the scalar complex-analytic
results proved here.

\section{Formalization ledger for version 1.0.0}\label{sec:formalization_ledger}

The preserved v1.0.0 release was built with Lean 4 v4.34.0-rc2 and Mathlib revision
\leanname{aeb72f1e9981b6d35b426b16dca7759be3e7ba06}
\cite{deMouraUllrich2021,mathlib2020}. Its 151 theorem-map rows resolve to
1164 unique declarations in the namespace \leanname{Komorebi.PaperIV}:
953 theorem declarations and 211 constructions. The recorded namespace audit scanned
2234 declarations, including 1999 theorem declarations. The historical warning fatal
build and supplied kernel audit allowed only
\leanname{propext}, \leanname{Classical.choice}, and
\leanname{Quot.sound} as standard transitive axioms.

\begingroup
\small
\begin{longtable}{@{}>{\raggedright\arraybackslash}p{0.23\textwidth}>{\raggedright\arraybackslash}p{0.70\textwidth}@{}}
\toprule
Manuscript records & Compiled Paper IV scope \\
\midrule
\endfirsthead
\toprule
Manuscript records & Compiled Paper IV scope (continued) \\
\midrule
\endhead
3.1--3.3 & Actual holomorphic root filters, convergent coefficient
selection, orthogonal resolution, the exact range--covariance characterization,
and the complete pointwise cyclic product law. The formal versions hold as well. \\
4.1--4.5 and 5.1--5.8 & Entire sections, finite exponential formulas,
derivative cycles, both Wronskians, and unique entire ODE basis expansions.
Both pointwise addition laws and denominator symmetries are checked.
The exact ordinary hypergeometric formula follows from the omitted-parameter
list, Gamma normalization, and Mathlib's everywhere convergent series. \\
6.1--6.6 & Quotient-algebra power bases and actual exponential-series
convergence in the finite-dimensional coefficient topology. Both pointwise
section-coordinate identities. Signed-shift powers, all matrix-exponential
entries, group laws and trace are checked. The characteristic polynomial
$X^N+1$, its distinct roots, exact spectrum and determinant one are proved. \\
7.1--7.8 & Named reciprocal families, exact parameter recurrences, Bell
formula, parameter degree and leading coefficient, and polynomial divisibility.
Parameter addition holds with any single fixed holomorphic logarithm on the
stated domain. Newton's recurrence gives the literal Hessenberg determinant.
All three compact reciprocal-polynomial formulas are checked. \\
8.1--8.5 & The full tangent coefficient series has the exact kappa and
factorial normalization, including $N=1$. Reciprocal differentiation holds on
every stated fixed logarithm domain. The actual higher Riccati identity holds
at every point where the denominator is nonzero. \\
9.1--9.10 & Arbitrary-parameter residue reciprocals and their analytic germs,
polynomial degree and leading term, sparse Appell laws, joint convolution,
monicity and covariance. The Hasse and derivative-truncation encodings are
proved equal. The named Appell polynomial satisfies both the forward action
and the exact $\Phi_{N,j}(\partial_x^N)^r$ inverse formula. \\
10.1--10.4 & Exact formal theta reconstruction, polynomial lift and degree,
literal divisibility by $u$ and, for even $N$, by $u+1$, and the theta Newton
recurrence with its first coefficient. \\
11.2, 11.4--11.11 & Two-sided maximum-modulus bounds prove exact order
and type. The indicator limsup formula holds in every direction, including ties.
Dominant-exponential errors and zero freeness are uniform on compact sets
avoiding ties. Near cancellation and the order-two zero set are checked.
For $N\geq3$, gap and remainder bounds, sector confinement, global exhaustion
and multiplicity identities give the full large-zero-string theorem.
The real leading correction, quadratic exponential error, signed derivative
asymptotic, eventual alternating sign and cubic correction are checked. \\
11.13--11.15 & The reduced series is entire, and
$\Dfun_N(z)=\Delta_N(z^N)$ preserves multiplicity away from zero. Finite
exceptional sets, eventual positive simple reduced zeros, and their stated
displacement are checked, including $N=2$. Both full disk-counting laws
retain all exceptional multiplicities. The full-zero reciprocal sums have
exact thresholds $s>1$ and $s>1/N$, including boundary divergence. \\
12.1 & For every $N\geq2$, the complete multiplicity-indexed product
converges uniformly on compact sets and equals $\Delta_N$. Jensen and Poisson
estimates prove rigidity of the actual quotient from derived growth bounds.
No growth certificate remains. The original-variable product for $\Dfun_N$
is also checked. \\
12.2--12.4 & The actual diagonal operator is trace class, with a
basis-independent trace and determinant given by limits of actual finite
compressions. Infinite complete homogeneous coefficients are limits of finite
monomial sums. Every normalized section has exact order $1/N$. For $N\geq2$
its full zero product and zero-moment Bell formulas are proved. The actual
complex-power germ at zero has the named residue Taylor coefficients. \\
12.6--12.10 & The normalized logarithm on each zero free disk has the
convergent zero-moment series. Its coefficients equal $\kappa_{N,m}$ and
satisfy Newton's recurrence. The Mittag--Leffler tangent expansion converges
absolutely and locally uniformly away from poles. Tangent coefficients and
the named reciprocal and theta families have their exact zero-moment Bell
formulas. The full tangent series is identified with the analytic function. \\
13.1--13.4 & Analytic division constructs each simple-zero principal part.
Its leading coefficient is uniquely determined. Removing the principal parts
gives a holomorphic remainder and the exact factorial-scaled Cauchy estimate.
Both leading-term errors are bounded, with the first absent for $r=1$.
Rotation gives the orbit-reduced main term. The secant formula holds as a
ratio limit, with an exponentially vanishing relative error. \\
14.2--14.4, 14.6 & The exact all-zero sums retain analytic multiplicity.
Their sharp absolute-convergence domains, boundary divergence and local
uniform convergence are checked. Both Hurwitz comparisons, entire corrections,
global meromorphic continuations, unique simple poles and exact residues are
proved. The order-two/four models have identically zero corrections. Branch
choices on the finite exceptional part change the continuation by an entire
function and preserve its pole and residue. \\
14.8 & The derivative-weighted series uses every actual simple zero and genuine
logarithms with bounded arguments. Absolute and local uniform convergence and
an entire sum are proved for $N\geq3$. For $N=2$, absolute/local uniform
convergence and holomorphy hold on $\Re s>1$. \\
\bottomrule
\end{longtable}
\endgroup

The declaration-level theorem map is authoritative for the preserved
v1.0.0 contract. Its release records full scoped coverage of all 68
theorem like records: 68 full, zero partial, and zero unmapped. Rows 143--151 close the final fifteen partial records
without altering their statements or assumptions. Imported standard fact
11.1 remains imported. Formal identities, normalized analytic germs, and
fixed logarithm domain identities retain their respective stated meanings.
Finite computations do not change these proof boundaries.

\section{Conclusion}\label{sec:conclusion}

Cyclic projection turns the exponential into a finite differential basis, and
the signed cycle produces the distinguished denominator $\Dfun_N$. Its
reciprocal powers yield Bell and Newton recurrences, an additive parameter
law, Appell systems, and theta polynomials. Reduction by $w=z^N$ then connects
the reciprocal coefficients to a genus zero product and reduced-zero power
sums.

Classical two term localization, specialized to the regular exponent polygon
and combined with reflection symmetry, places all sufficiently large zeros
on symmetry rays. The displacement, derivative and counting estimates yield
the family-specific Hurwitz comparison for the orbit zero zeta.

\begingroup
\small
\let\prscbibliography\thebibliography
\renewcommand{\thebibliography}[1]{\prscbibliography{#1}\setlength{\itemsep}{0pt}\setlength{\parsep}{0pt}}
\bibliographystyle{alpha}
\bibliography{references}
\endgroup

\clearpage
\appendix
\section{Visual atlas}\label{app:visual-atlas}
The twenty-six plates retain the worked models of the separate visual guide. Each stable label V1 through V26 names a mathematical mechanism and points back to the relevant statements. The plates distinguish exact identities, rational certificates, sampled curves, and inherited formal evidence.
\begin{readerroadmap}[title={Three routes through the atlas}]
For algebra, read V3 through V16. For zero geometry and its consequences, read V17 through V24. V25 follows one coefficient through the complete calculation. V26 explains the evidence attached to the models.
\end{readerroadmap}
\clearpage
\begingroup
\fontsize{10}{12}\selectfont\setlength{\parskip}{4pt}
\setlength{\abovedisplayskip}{6pt}\setlength{\belowdisplayskip}{6pt}
\vpage{1}{A coefficient can become zero data}{READER'S GUIDE}{Traditional monograph: Section 2 (p. 2), Section 7 (p. 13), Section 11 (p. 25), Section 12 (p. 34).}
This guide follows a concrete analytic family through three descriptions: selected Taylor coefficients, a finite differential system, and the zeros of an entire function. The formulas are linked, but the evidence needed for each description is different.
\vfig{three_views}
\head{Keep three small models in view}
For $N=2$, the signed denominator is $D_2(z)=\cos z$. Its reciprocal is the familiar secant series. For $N=3$, a third exponential perturbs the zeros of a dominant cosine. For $N=4$, an exact factorization gives a complete zero picture:
\[
D_4(z)=\cos(z/\sqrt2)\cosh(z/\sqrt2).
\]
The examples move between these orders deliberately. Order two makes the coefficient arithmetic transparent. Order three shows a genuine displacement. Order four makes orbit reduction completely explicit.
\head{Three routes through the guide}
\par\noindent\begin{tabularx}{\linewidth}{@{}lY@{}}\toprule
Route & Pages and question\\\midrule
Algebra first & \pagerefV{3} to \pagerefV{16}: how do selection, signs, reciprocals, and finite differential operators work?\\
Zeros first & \pagerefV{17} to \pagerefV{24}: where can cancellation happen, and how does it determine coefficients?\\
One full calculation & \pagerefV{25}: derive the same number, $69$, from denominator coefficients and from zero moments.\\\bottomrule
\end{tabularx}
\takeaway{The pictures carry computable objects. They do not replace the global analytic arguments or extend the inherited Lean contract.}
\platecaption{Exact identities and rational certificates are distinguished from sampled curves. All plates are reproducible vector drawings. The source package preserves the original proofread manuscript.}

\vpage{2}{Two variables, several normalizations}{OBJECT ATLAS}{Traditional monograph: Section 2 (p. 2), Definition 5.7 (p. 9), Definition 7.1 (p. 13), Definition 9.1 (p. 19).}
The most important coordinate change is $w=z^N$. A term of degree $n$ in $w$ has degree $Nn$ in $z$. Its coefficient does not change, but an exponential-generating-function label does.
\[
\underbrace{b_{N,n}(r)\,w^n}_{\text{ordinary coefficient}}
=\underbrace{P_{N,n}(r)\frac{z^{Nn}}{(Nn)!}}_{\text{block exponential coefficient}},
\qquad P_{N,n}(r)=(Nn)!b_{N,n}(r).
\]
\par\noindent\begin{tabularx}{\linewidth}{@{}lY@{}}\toprule
Object & Definition and role\\\midrule
$E_{N,j}(z)$ & $\displaystyle\sum_{n\ge0}\frac{z^{Nn+j}}{(Nn+j)!}$ selects one residue class of the exponential.\\
$D_{N,j}(z)$ & The same series with $(-1)^n$. One full derivative cycle contributes a minus sign.\\
$D_N(z)$ & $D_{N,0}(z)$, the denominator whose zeros control the signed reciprocal family.\\
$\Delta_N(w)$ & $\displaystyle\sum_{n\ge0}\frac{(-1)^nw^n}{(Nn)!}$. Thus $D_N(z)=\Delta_N(z^N)$.\\
$P_{N,n}(r)$ & Block-EGF coefficients of $D_N^{-r}$. They are polynomials in $r$.\\
$\kappa_{N,m}$ & $\displaystyle-\log\Delta_N(w)=\sum_{m\ge1}\kappa_{N,m}w^m/m$.\\
$T_N(z)$ & $-D_N'(z)/D_N(z)$, a function, not the coefficient $T_N(n)$.\\
$\Phi_{N,j}(w)$ & $\displaystyle j!\sum_{n\ge0}\frac{w^n}{(Nn+j)!}$ removes the forced zero of $E_{N,j}$.\\
$C_{N,j}(z)$ & $1/\Phi_{N,j}(z^N)$. This is generally not $1/D_N(z)$.\\
$\Lambda_N$ & Zeros of $\Delta_N$, counted with multiplicity, one per rotational orbit of zeros of $D_N$.\\\bottomrule
\end{tabularx}
\head{A normalization test}
The coefficient $5/24$ of $w^2$ in $\Delta_2(w)^{-1}$ is the same term as $5z^4/4!$. It is not $5w^2/2!$. Keeping the two factorials separate prevents errors in Bell formulas, convolution weights, and zero-moment reconstruction.
\takeaway{Use $z$ for differentiation and rotational geometry. Use $w$ for reciprocal recurrences, canonical products, and orbit-zero sums.}

\vpage{3}{A root filter is an exact vote}{CYCLIC PROJECTION}{Traditional monograph: Theorem 3.1 (p. 4), Proposition 3.2 (p. 5).}
For $N=4$, write $\omega=i$. To extract residue $j$, evaluate the function at four rotated arguments and attach the phase $i^{-j\ell}$:
\[
(\mathcal B_{4,j}F)(z)=\frac14\sum_{\ell=0}^{3}i^{-j\ell}F(i^\ell z).
\]
A monomial $a_mz^m$ is multiplied by the finite average of $i^{(m-j)\ell}$. Matching phases all vote $1$. Nonmatching phases cancel exactly.
\vfig{root_filter}
\head{An example that can be checked coefficient by coefficient}
Take $F(z)=\sum_{m=0}^{12}(m+1)z^m$. Then
\[
\mathcal B_{4,2}F=3z^2+7z^6+11z^{10}.
\]
Applying the same selector again changes nothing. Applying a different selector gives zero. Summing the four selected polynomials returns the original $F$.
\head{Selection is also a symmetry condition}
For the selected polynomial $F_2$, rotation gives $F_2(iz)=-F_2(z)$. More generally, a residue-$j$ section satisfies $F_j(\omega z)=\omega^jF_j(z)$. The phase label therefore describes both the coefficient addresses and the response to rotation.
\takeaway{This is not an approximate Fourier fit. The finite root sum is exactly $0$ or $1$ on every coefficient.}
\platecaption{The checker compares the actual complex average with independent coefficient selection for this polynomial. The drawing shows the four complex phases, not a numerical convergence argument.}

\vpage{4}{Products mix the residue lanes}{CYCLIC CONVOLUTION}{Traditional monograph: Theorem 3.3 (p. 5).}
Selection separates a function additively. It does not make multiplication diagonal. A residue-$i$ term multiplied by a residue-$j$ term lands in residue $i+j\pmod N$.
For $N=3$, use
\[
F=1+2z+3z^2+z^3,\qquad G=2+z+z^2.
\]
The table displays every product of their three selected sections.
\vfig{cyclic_product}
The shaded cells contribute to residue zero:
\[
2(1+z^3)+(2z)(z^2)+(3z^2)z=2+7z^3.
\]
Keeping only the product of the residue-zero inputs would give $2+2z^3$ and lose five units of the $z^3$ coefficient. The mixed cells are part of the exact result.
\head{The general rule}
\[
\mathcal B_{N,k}(FG)=
\sum_{i+j\equiv k\ (N)}(\mathcal B_{N,i}F)(\mathcal B_{N,j}G).
\]
Every pair belongs to exactly one output residue. No term is discarded and no coefficient is estimated. This finite cyclic convolution is the algebra behind the later addition formulas and block-coefficient recurrences.
\takeaway{A packet label tells you where an interaction lands, not that interactions with other packets vanish.}
\platecaption{All nine table entries and all three output sections are checked as exact polynomial identities. This is the scalar cyclic specialization used by the paper.}

\vpage{5}{One signed edge determines the cycle}{DIFFERENTIAL SYSTEM}{Traditional monograph: Theorem 5.3 (p. 8), Theorem 6.4 (p. 11), Corollary 6.5 (p. 12), Corollary 6.6 (p. 12).}
Let $J_3e_0=e_1$, $J_3e_1=e_2$, and $J_3e_2=-e_0$. A complete turn changes sign, so $J_3^3=-I$. The signed sections are the coordinates of $e^{zJ_3}e_0$.
\vfig{signed_shift}
Here $D_j$ inside the matrix means $D_{3,j}(z)$. For example,
\[
D_{3,0}=1-\frac{z^3}{3!}+\frac{z^6}{6!}-\cdots,
\qquad D_{3,1}=z-\frac{z^4}{4!}+\cdots.
\]
Rows are output coordinates and columns are input coordinates. The upper-triangular minus signs record crossing the signed edge while moving from the column to the row. They are not arbitrary orientation choices.
\head{Two checks connect the matrix and the scalar equation}
The initial derivative matrix is the identity. Each initial jet $(y(0),y'(0),y''(0))$ therefore determines a unique linear combination. Also,
\[
\frac13\Tr(e^{zJ_3})=D_3(z),\qquad \det(XI-J_3)=X^3+1.
\]
The roots of $-1$ appearing later in the exponential sum are already the spectrum of this concrete matrix.
\takeaway{The functions, their initial data, the addition law, and the exponent polygon come from one signed shift.}
\platecaption{The exact tests check every matrix-entry Taylor coefficient through degree 14 for $N=1,\ldots,5$. The displayed determinant-one law is for $N\ge2$. $N=1$ is a separate boundary case.}

\vpage{6}{Addition is quotient multiplication}{FINITE ALGEBRA}{Traditional monograph: Proposition 5.5 (p. 8), Theorem 6.2 (p. 10).}
Work in $\CC[\xi]/(\xi^3+1)$. Every element has coordinates in $1,\xi,\xi^2$. The entire exponential has the finite-coordinate form
\[
e^{z\xi}=D_{3,0}(z)+D_{3,1}(z)\xi+D_{3,2}(z)\xi^2.
\]
The multiplication table makes the carry sign unavoidable.
\vfig{signed_addition}
Multiply the coordinate forms of $e^{z\xi}$ and $e^{w\xi}$. Their product is $e^{(z+w)\xi}$. The coefficient of $1$ receives one contribution without a wrap and two contributions with a wrap. Reading that coefficient gives the displayed addition law.
\head{Compare the unsigned model}
In $\CC[\xi]/(\xi^3-1)$, the same table has no minus signs. Its exponential coordinates are the unsigned sections $E_{3,j}$. Thus the difference between the two differential systems is a single algebraic relation, not a change in how coefficients are grouped.
For general $N$, the carry factor is
\[
(-1)^{\lfloor(i+j)/N\rfloor},\qquad 0\le i,j<N.
\]
Only zero or one wrap is possible in a product of two coordinate basis elements.
\takeaway{The twisted addition law is the ordinary exponential group law expressed in a signed cyclic algebra.}
\platecaption{Bivariate Taylor identities are checked through total degree ten in every section for $N=2,3,4$. Omitting the carry sign is one of the deliberately rejected controls.}

\vpage{7}{Why a hypergeometric name appears}{FACTORIAL NORMALIZATION}{Traditional monograph: Theorem 4.5 (p. 7).}
The hypergeometric form is another way to organize the same coefficients. It does not introduce an extra function into the construction. Split a factorial into its residue classes and divide each factor by $N$:
\[
(Nn+j)!=j!N^{Nn}n!\prod_{a\in A_{N,j}}(a)_n.
\]
The parameter equal to $1$ is omitted from $A_{N,j}$ because the hypergeometric convention already supplies $(1)_n=n!$.
\vfig{hypergeometric}
\head{Check the displayed case}
For $N=3,j=1,n=2$, the six factors after the initial $1!$ are $2,3,4,5,6,7$. Dividing by three and grouping them gives the three rows of the plate. The middle row supplies $2!$. The others supply $(2/3)_2$ and $(4/3)_2$. The right-hand side is exactly $5040$.
Consequently,
\[
E_{3,1}(z)=z\left(1+\frac{z^3}{4!}+\frac{z^6}{7!}+\cdots\right).
\]
The prefactor $z$ retains the forced zero at the origin. The hypergeometric argument $z^3/27$ retains the scaling $N^N$.
\takeaway{Parameters, the omitted unit parameter, the prefactor, and the argument scaling all come from factorial bookkeeping.}
\platecaption{Five choices of $(N,j)$ are checked against the exact factorial identity for $n=0,\ldots,7$. These finite checks support the example. The general analytic identification belongs to the inherited manuscript.}

\vpage{8}{Four zeros become one orbit zero}{THE REDUCED VARIABLE}{Traditional monograph: Definition 5.7 (p. 9), Example 11.12 (p. 31), Corollary 11.13 (p. 32).}
For $N=4$ the factorization is exact:
\[
D_4(z)=\cos(z/\sqrt2)\cosh(z/\sqrt2),\qquad D_4(z)=\Delta_4(z^4).
\]
Set $A_4=\sqrt2\pi$. All the zeros are $A_4(m+1/2)i^k$, with $m\ge0$ and $0\le k<4$. Each zero is simple.
\vfig{orbit_reduction}
The drawing uses $u=z/A_4$ and $v=w/A_4^4$, so its marked coordinates are rational. The first four-point orbit has fourth power $1/16$. The second has fourth power $81/16$.
\head{Multiplicity is not multiplied in the reduced function}
At a nonzero point, $z\mapsto z^4$ has nonzero derivative. It preserves the local multiplicity of a zero. Four distinct simple zeros of $D_4$ therefore correspond to one simple zero of $\Delta_4$, not one zero of multiplicity four.
The reduced zeros are
\[
\lambda_m=A_4^4(m+1/2)^4=4\pi^4(m+1/2)^4.
\]
This is why an orbit factor in the $z$ plane is $1-z^4/\lambda_m$. It is the product of the four corresponding linear factors, with their phases included.
\takeaway{Orbit reduction removes repetition. It changes growth and counting exponents, while preserving each nonzero zero's local multiplicity.}
\platecaption{The quartic factorization, its differential equation, and the displayed orbit powers are checked independently. Here the complete zero set is known exactly. A finite numerical root search is not being used to establish it.}

\vpage{9}{A reciprocal is a triangular solve}{COEFFICIENT RECONSTRUCTION}{Traditional monograph: Theorem 7.3 (p. 14), Theorem 7.6 (p. 15).}
Start in the reduced variable with
\[
\Delta_2(w)=1-\frac w2+\frac{w^2}{24}-\frac{w^3}{720}+\cdots.
\]
Write $1/\Delta_2(w)=\sum b_nw^n$. Matching the product with $1$ gives a unit lower-triangular system. At each row there is exactly one new unknown.
\vfig{triangular_inverse}
The first three equations are
\[
b_1-\tfrac12=0,\qquad b_2-\tfrac12b_1+\tfrac1{24}=0,
\qquad b_3-\tfrac12b_2+\tfrac1{24}b_1-\tfrac1{720}=0.
\]
They yield $b_1=1/2$, $b_2=5/24$, and $b_3=61/720$. Multiplication by $(2n)!$ produces the secant coefficients $1,5,61$.
\head{Why this is an exact finite calculation}
To compute the coefficient of $w^n$, only denominator coefficients through $w^n$ are needed. Later terms of the entire series cannot affect it. A finite series inverse is therefore exact at its declared truncation order, even though evaluating that truncated series is only an approximation to the full analytic function away from the origin.
\takeaway{A unit constant term makes coefficient recovery triangular. It does not by itself specify the radius of analytic convergence.}
\platecaption{The checks compare an independent binomial-series expansion, the parameter recurrence, and the logarithmic Newton recurrence through block index six. The table uses ordinary coefficients, with factorial conversion shown separately.}

\vpage{10}{The reciprocal order is a polynomial}{BELL PARTITIONS}{Traditional monograph: Theorem 7.2 (p. 14), Corollary 7.7 (p. 16).}
Let $r$ vary instead of fixing the reciprocal power. Since
\[
\Delta_N(w)^{-r}=\exp\!\left(r\sum_{m\ge1}\kappa_{N,m}\frac{w^m}{m}\right),
\]
the coefficient of $w^n$ is a polynomial of degree $n$ in $r$. Bell polynomials collect the possible ways logarithmic terms contribute total degree $n$.
\vfig{bell_partitions}
For $N=2$, the first three logarithmic coefficients are
\[
\kappa_1=\frac12,\qquad\kappa_2=\frac16,\qquad\kappa_3=\frac1{15}.
\]
The third Bell polynomial is $B_3(x_1,x_2,x_3)=x_1^3+3x_1x_2+x_3$. Substituting $x_1=r\kappa_1$, $x_2=r\kappa_2$, $x_3=2r\kappa_3$ and multiplying by $6!/3!$ gives
\[
P_{2,3}(r)=15r^3+30r^2+16r.
\]
At $r=1$ this is $61$. At $r=2$ it is $272$. At $r=0$ it vanishes, as every positive-index reciprocal coefficient must.
\head{The leading term has a simple source}
Only $n$ copies of the linear logarithmic term can contribute $r^n w^n$. Therefore
\[
[r^n]P_{N,n}(r)=\frac{(Nn)!}{n!(N!)^n}.
\]
\takeaway{The polynomial parameter records how many logarithmic pieces are combined. It is not a numerical interpolation in the reciprocal order.}
\platecaption{The partition plate is a finite combinatorial expansion. The exact tests verify the leading coefficient and parameter polynomial through block index six for five orders.}

\vpage{11}{Two recurrences reach the same answer}{NEWTON AND ADDITION}{Traditional monograph: Theorem 7.3 (p. 14), Theorem 7.4 (p. 15), Theorem 7.6 (p. 15).}
One recurrence uses the denominator coefficients. Another uses its logarithm. With $b_n(r)=P_{N,n}(r)/(Nn)!$, the shorter logarithmic form is
\[
nb_n(r)=r\sum_{k=1}^{n}\kappa_{N,k}b_{n-k}(r),\qquad b_0=1.
\]
The negative superdiagonal in the Newton determinant is essential.
\vfig{newton_addition}
\head{The additive parameter law has block weights}
Using the same logarithm, $D_N^{-r}D_N^{-s}=D_N^{-(r+s)}$. Thus
\[
P_{N,n}(r+s)=\sum_{q=0}^{n}\binom{Nn}{Nq}P_{N,q}(r)P_{N,n-q}(s).
\]
For $N=2,n=3,r=s=1$, the four weights are $1,15,15,1$. The two middle terms are $15\cdot1\cdot5$ and $15\cdot5\cdot1$. Together with the endpoints, they give $272$.
Using $\binom nq$ instead would be incorrect for these coefficients. That ordinary binomial law belongs to the separately normalized sequence $n!P_{N,n}/(Nn)!$.
\takeaway{Agreement of independent recurrences is a useful check precisely because their bookkeeping is different. Neither permits silently changing the factorial convention.}
\platecaption{The checker compares both recurrence routes and the determinant, then checks the parameter-addition identity symbolically. A deliberately incorrect ordinary-binomial substitution is rejected.}

\vpage{12}{Compress the ODE into a tangent}{LOGARITHMIC DIFFERENTIATION}{Traditional monograph: Proposition 8.2 (p. 18), Corollary 8.4 (p. 18), Theorem 8.5 (p. 19).}
The transfer tangent is $T_N=-D_N'/D_N$. It compresses the linear differential equation into a nonlinear scalar identity. This division is allowed only on the zero-free set of the denominator.
\vfig{riccati}
The coefficients retain the reduced logarithmic information:
\[
T_N(z)=N\sum_{m\ge1}\kappa_{N,m}z^{Nm-1},
\qquad T_N(n)=N(Nn+N-1)!\kappa_{N,n+1}.
\]
For order four, $\kappa_{4,1}=1/24$ and $\kappa_{4,2}=17/10080$. These give $T_4(0)=1$ and $T_4(1)=34$, explaining the last line of the plate.
\head{A nonlinear equation with a precise origin}
If $D=e^g$ locally, then $g'=-T$. Differentiating $e^g$ three times gives $D'''/D=-T''+3TT'-T^3$. Setting $D'''/D=-1$ yields the cubic Riccati identity shown above. No additional nonlinear model has been assumed.
The reciprocal family also satisfies $(D_N^{-r})'=rT_ND_N^{-r}$ on a fixed-log domain. At a zero of $D_N$, these meromorphic expressions have poles rather than finite values.
\takeaway{The nonlinear tangent equation is exact where division is legal. It is not a global regular differential equation across the zeros.}
\platecaption{The tests verify the tangent coefficient recurrence, the exact cubic ODE, and the symbolic Riccati reduction after clearing denominators. The order-one tangent is the separate constant function $T_1=1$.}

\vpage{13}{A complex power needs a domain}{BRANCH BOOKKEEPING}{Traditional monograph: Section 2 (p. 2), Definition 7.1 (p. 13), Corollary 8.4 (p. 18).}
A formal unit series has a formal logarithm at the origin. To interpret that germ analytically, choose a simply connected zero-free region containing $0$ and normalize the logarithm to vanish there.
\[
D_N(z)^{-r}=\exp(-rL(z)),\qquad e^{L(z)}=D_N(z),\quad L(0)=0.
\]
\vfig{branch_domain}
For $D_2=\cos z$, the open disk $|z|<\pi/2$ is an admissible domain. A small loop around the simple zero $\rho=\pi/2$ is not contained in that disk. Analytic continuation around the zero changes the local logarithm by $2\pi i$, and the power by $e^{-2\pi ir}$.
\head{Three different kinds of statement}
A \emph{formal} identity compares coefficients at the origin. A \emph{holomorphic} identity for complex powers uses a common specified logarithm. A positive \emph{integer} reciprocal power is globally meromorphic and needs no branch choice. These scopes cannot be exchanged without a further argument.
For example, $D^{-r}D^{-s}=D^{-(r+s)}$ is immediate using one fixed logarithm. Combining values built from unrelated branches can introduce additional phase factors.
\takeaway{The branch choice is mathematical data. It is not a plotting convention that can be removed by drawing a cleaner curve.}
\platecaption{The domain picture is schematic. The phase change comes from the local simple-zero model and winding number, not from a sampled loop. Later nearest-pole asymptotics retain positive integral reciprocal order.}

\vpage{14}{Three similar reciprocals differ}{NORMALIZED SECTIONS}{Traditional monograph: Definition 9.1 (p. 19), Theorem 9.2 (p. 20), Example 9.11 (p. 22).}
The signed denominator and the normalized unsigned sections produce different reciprocal families. In order two,
\[
D_2^{-1}=\sec z,\qquad C_{2,0}=\sech z,\qquad C_{2,1}=z/\sinh z.
\]
Their first block-EGF coefficients make the difference visible.
\vfig{different_reciprocals}
\head{Remove the forced zero before taking a reciprocal}
The section $E_{N,j}$ begins with $z^j/j!$. Define
\[
E_{N,j}(z)=\frac{z^j}{j!}\Phi_{N,j}(z^N),\qquad\Phi_{N,j}(0)=1.
\]
Then $C_{N,j}(z)=1/\Phi_{N,j}(z^N)$ has value $1$ at the origin. For $j>0$, the notation $(z^j/j!)/E_{N,j}(z)$ means this removable continuation. It does not mean evaluating an undefined $0/0$ expression.
For $N=2,j=1$, the first nontrivial coefficient is $-1/3$ in front of $z^2/2!$, equivalently $-z^2/6$. Dropping the numerator $z$ would instead produce the pole of $1/\sinh z$.
\takeaway{A minus sign in the denominator series and removal of an initial zero are different operations. Their reciprocal coefficients should not be interchanged.}
\platecaption{All entries in the table are exact rationals computed from independent triangular inverses. The table explicitly uses block-EGF columns. Ordinary coefficients are obtained by dividing by the displayed factorials.}

\vpage{15}{An infinite symbol acts finitely}{APPELL POLYNOMIALS}{Traditional monograph: Theorem 9.6 (p. 21), Theorem 9.8 (p. 21), Proposition 9.9 (p. 21).}
For $C_{2,0}(z)=\sech z$, define $E_m(x)$ by $\sech z\,e^{xz}=\sum E_m(x)z^m/m!$. The formal differential-operator formula is
\[
E_m(x)=\sech(\partial_x)x^m.
\]
Although the symbol is an infinite series, derivatives above degree $m$ vanish.
\vfig{appell_operator}
For $m=4$, only three terms survive:
\[
\left(1-\frac{\partial_x^2}{2!}+5\frac{\partial_x^4}{4!}\right)x^4
=x^4-6x^2+5.
\]
The four real roots of this polynomial are $\pm1,\pm\sqrt5$. They are not the zeros of the original analytic denominator. The plot distinguishes the input polynomial and its transformed output.
\head{An inverse operation, not a truncation error}
Applying $\cosh(\partial_x)$ returns $x^4$ exactly. The general inverse uses the reciprocal formal symbol. The lowering and translation identities are
\[
E_m'=mE_{m-1},\qquad E_m(x+y)=\sum_{q=0}^m\binom mqE_q(x)y^{m-q}.
\]
Both follow from the generating function. No operator-series convergence assumption on an infinite-dimensional function space is needed for this polynomial calculation.
\takeaway{Formal substitution into the derivative is safe here because every polynomial sees only finitely many derivatives. That is a domain statement.}
\platecaption{The polynomial, its inverse image, and the Appell identities are checked symbolically through degree eight in four normalized-section families. The displayed curves are samples of exact polynomials.}

\vpage{16}{Scale and order enter differently}{THETA COUPLING}{Traditional monograph: Theorem 10.2 (p. 23), Proposition 10.3 (p. 23), Proposition 10.4 (p. 23).}
Let $H_N(z)=E_{N,0}(z)$. The theta family uses the same parameter $u$ twice, but in two different roles:
\[
H_N(uz)^{-1}H_N(z)^{-u}
=\sum_{n\ge0}\Theta_{N,n}(u)\frac{z^{Nn}}{(Nn)!}.
\]
Argument scaling contributes $u^{Nm}$ at logarithmic block $m$. Reciprocal order contributes $u$.
\vfig{theta}
Their combined logarithm has coefficient
\[
(-1)^m\kappa_{N,m}\frac{u^{Nm}+u}{m}.
\]
At the first block this gives $\Theta_{N,1}(u)=-(u^N+u)$. Every positive-index theta polynomial vanishes at $u=0$. For even $N$, $H_N(-z)=H_N(z)$, so the defining product is also $1$ at $u=-1$.
\head{The odd-order boundary is visible}
For $N=3$, $\Theta_{3,1}(-1)=2$, not zero. The factor $u+1$ is not a uniform law across orders. In contrast,
\[
\Theta_{2,2}(u)=u(u+1)(5u^2+u+2).
\]
The two guaranteed factors follow from the whole generating function, not from guessing the first few polynomial roots.
\takeaway{A common parameter does not make argument scaling and reciprocal order interchangeable. Their powers must be tracked separately.}
\platecaption{The exact checker compares theta convolution with its logarithmic Newton recurrence through block index four. The odd-order factor claim is explicitly rejected as an invalid control.}

\vpage{17}{Growth comes from a polygon}{DOMINANT EXPONENTS}{Traditional monograph: Definition 11.3 (p. 26), Theorem 11.4 (p. 26).}
Write the signed denominator as a finite exponential sum:
\[
D_N(z)=\frac1N\sum_{\eta^N=-1}e^{\eta z}.
\]
For $z=re^{i\theta}$, the size of one exponential is $e^{r\Re(\eta e^{i\theta})}$. The maximal projection determines the directional exponential growth.
\vfig{indicator}
The indicator is
\[
h_N(\theta)=\max_{0\le k<N}\cos\left(\theta+\frac{(2k+1)\pi}{N}\right).
\]
Away from a tie direction, one exponential dominates all the others by a uniform positive projection gap on compact angular sets. For sufficiently large radius, smaller terms cannot cancel it. Those sectors are eventually zero-free.
\head{Why zeros remain possible on a growing ray}
At a tie, two adjacent exponents have the same real projection. Their complex phases can be opposite. Their sum can therefore vanish even though each term is exponentially large. The minima in the order-four indicator plot occur precisely at the tie directions $\theta=k\pi/2$.
The indicator uses a limsup in radius. It does not assert that $r^{-1}\log|D_N(re^{i\theta})|$ has an ordinary limit through the zeros on a tie ray.
\takeaway{The polygon predicts where large cancellation is possible. It does not, by itself, enumerate zeros or prove their simplicity.}
\platecaption{The order-four polygon and support-function formula are explicit. The curve is a sampled display of that formula. Eventual localization and uniqueness require the separate zero-analysis arguments.}

\vpage{18}{Two large terms nearly cancel}{THE CUBIC MODEL}{Traditional monograph: Lemma 11.5 (p. 27), Theorem 11.7 (p. 29), Corollary 11.9 (p. 31).}
For real $t$, order three has the exact formula
\[
D_3(t)=\frac{e^{-t}+2e^{t/2}\cos(\sqrt3t/2)}{3}.
\]
Divide by the nonzero factor $2e^{t/2}/3$. The resulting equation separates the dominant cosine from the small third exponential.
\vfig{two_term}
Ignoring the last term gives model zeros
\[
t_m=\frac{2\pi}{\sqrt3}(m+1/2).
\]
At $t_m$, the full normalized expression equals $\tfrac12e^{-3t_m/2}>0$, so $t_m$ is not an exact zero. Linearizing the cosine at that point predicts
\[
\delta_m^{(1)}=\frac{(-1)^m}{\sqrt3}e^{-3t_m/2}.
\]
The alternating slope explains why the true zero is to the right of the first model zero, then to the left of the second, and so on in the eventual regime.
\head{Symmetry needs uniqueness}
Real coefficients alone allow nonreal conjugate pairs. The manuscript first proves that a near-model disk contains exactly one zero. Conjugation preserves that disk, so the unique zero must equal its conjugate. This additional uniqueness step, not the symmetry alone, places it exactly on the ray.
\takeaway{The correction is caused by a real, nonzero remainder. A good cosine approximation does not make its model zeros exact.}
\platecaption{The curves are sampled analytic expressions. The next page independently certifies four narrow real-zero intervals rather than trusting the plotted crossing or the floating-point root finder.}

\vpage{19}{A root estimate becomes a certificate}{EXACT RATIONAL BRACKETS}{Traditional monograph: Theorem 11.7 (p. 29), Theorem 11.8 (p. 30), Corollary 11.9 (p. 31).}
The root finder is used only to suggest intervals. Each displayed interval is accepted by exact rational arithmetic: opposite endpoint signs prove existence, and an interval bound for $D_3'$ excludes zero throughout the interval, proving uniqueness inside it.
\vfig{cubic_certificates}
\head{How the certificate works}
Evaluate the Taylor polynomial
\[
S_{90}(t)=\sum_{n=0}^{90}\frac{(-1)^nt^{3n}}{(3n)!}
\]
at the two rational endpoints. From the next term onward, magnitudes decrease, because their successive ratio is
\[
\frac{t^3}{(3n+1)(3n+2)(3n+3)}<1
\]
there. The alternating-series remainder is bounded by the first omitted term. A separate outward interval evaluation of the differentiated polynomial, plus its tail bound, certifies the derivative sign.
The first bracket is
\[
1.84981279919014347629<\tau_0<1.84981279919014347630.
\]
Its displacement differs from the leading asymptotic prediction by about five percent. The later examples have much smaller relative discrepancies, as the plate records.
\takeaway{A small residual is not a root certificate. Endpoint signs plus a derivative bound give a checkable local existence-and-uniqueness statement.}
\platecaption{Four certified intervals do not establish that no other real or complex zeros exist. The displacement columns are high-precision diagnostics, not a finite verification of a global big-$O$ theorem. All rational enclosures are included in the checks file.}

\vpage{20}{Counting changes under the power map}{EXACT FOURTH-ORDER STRINGS}{Traditional monograph: Example 11.12 (p. 31), Theorem 11.14 (p. 32).}
Order four gives a case with no asymptotic displacement at all. In normalized coordinates the positive radial locations are $u=m+1/2$, while the reduced-zero locations are $v=(m+1/2)^4$.
\vfig{counting}
The two rows describe the same zero orbits. The upper row counts four zeros each time a new radius is reached. The lower row counts one reduced zero each time its fourth-power location is reached. Both formulas include zeros on the boundary of the disk.
\head{Separate radius conversion from orbit multiplicity}
In unscaled variables the general relation is
\[
n_{D_N}(R^{1/N})=Nn_{\Delta_N}(R).
\]
There are two operations here: replace the radius by its $N$th root and divide the count by $N$. Forgetting either operation changes the leading constant.
For this family the asymptotics are
\[
n_{D_N}(R)=\frac{N\sin(\pi/N)}\pi R+O(1),\qquad
n_{\Delta_N}(R)=\frac{\sin(\pi/N)}\pi R^{1/N}+O(1).
\]
The different powers agree with the entire-function orders $1$ and $1/N$.
\takeaway{The reduced function has fewer repeated zeros and a different radius scale. Its counting law is not obtained by dividing a same-radius count alone.}
\platecaption{The exact order-four staircase formulas are tested at rational radii including jump points. The large-zero theorem for general $N$ permits finitely many exceptions. These exact order-four identities do not justify discarding that boundary elsewhere.}

\vpage{21}{A product and an operator share zeros}{CANONICAL PRODUCT}{Traditional monograph: Theorem 12.1 (p. 34), Theorem 12.2 (p. 35).}
For $N\ge2$, the reduced entire function has order $1/N<1$. Its multiplicity-sensitive genus-zero product is
\[
\Delta_N(w)=\prod_{\lambda\in\Lambda_N}(1-w/\lambda).
\]
Putting $K_Ne_\lambda=\lambda^{-1}e_\lambda$ on $\ell^2(\Lambda_N)$ gives the corresponding diagonal Fredholm determinant. Repeated zeros require repeated diagonal entries.
\vfig{product}
\head{Finite compression is an approximation}
For $N=2$, $\lambda_m=\pi^2(m+1/2)^2$. The finite determinant $P_L$ is exactly the displayed finite product, but it is not exactly $\Delta_2$. For $0\le w\le1$ and $L\ge1$, an elementary tail estimate gives
\[
0\le1-\frac{\Delta_2(w)}{P_L(w)}
\le w\sum_{m\ge L}\lambda_m^{-1}
\le\frac1{9(L-1/2)}.
\]
The last bound uses $\pi>3$ and integral comparison. Thus truncation has an explicit error mechanism.
\takeaway{The infinite product is justified by the analytic theorem. Finite determinants provide controlled approximations, not substitutes for that theorem.}
\platecaption{The exact finite determinant test includes a repeated eigenvalue and rejects deduplication. The curves use known cosine zeros. Order one is excluded: $\Delta_1(w)=e^{-w}$ has no zeros, so its empty product would be $1$, not the function.}

\vpage{22}{Finite coefficients give infinite moments}{ZERO POWER SUMS}{Traditional monograph: Theorem 12.6 (p. 36), Corollary 12.7 (p. 37), Corollary 12.9 (p. 38), Corollary 12.10 (p. 38).}
Taking the logarithm of the canonical product identifies
\[
\kappa_{N,q}=\sigma_{N,q}:=\sum_{\lambda\in\Lambda_N}\lambda^{-q}
=\Tr(K_N^q),\qquad q\ge1.
\]
The same exact numbers can be computed from finitely many denominator coefficients by Newton's recurrence. No numerical zero list is needed for that calculation.
\vfig{moments}
For order two,
\[
\sigma_{2,1}=\frac12,\quad\sigma_{2,2}=\frac16,\quad
\sigma_{2,3}=\frac1{15},\quad\sigma_{2,4}=\frac{17}{630}.
\]
The plate shows how successive inverse powers emphasize the nearest zero. Its shaded pieces are contributions from the first five positive reduced zeros, followed by the remaining tail. The total in each row is the exact rational moment.
\head{The tangent is a sum of poles}
Logarithmic differentiation also gives
\[
T_N(z)=Nz^{N-1}\sum_{\lambda\in\Lambda_N}\frac1{\lambda-z^N}
\]
away from poles. Summability of $|\lambda|^{-1}$ permits the locally uniform partial-fraction expansion without added regularizing terms.
\takeaway{A global analytic product theorem connects the infinite zero sum to a finite algebraic recurrence. The recurrence alone does not prove the product theorem.}
\platecaption{Displayed moments are checked against formal logarithm coefficients and exact order-two/order-four zeta values. Bar widths are numerical shares. General power sums use $\lambda^{-q}$, not $|\lambda|^{-q}$ when exceptional nonpositive zeros are present.}

\vpage{23}{Nearest poles set coefficient growth}{SINGULARITY ANALYSIS}{Traditional monograph: Theorem 13.1 (p. 39), Corollary 13.2 (p. 39), Corollary 13.3 (p. 39), Corollary 13.4 (p. 40).}
For $1/\cos z$, the nearest poles are $\rho=\pm\pi/2$. They contribute equally to even-degree coefficients. Including both poles gives
\[
P_{2,n}(1)\sim2(2n)!\left(\frac2\pi\right)^{2n+1}.
\]
The factor two is orbit multiplicity, not an adjustable fitted constant.
\vfig{pole_asymptotic}
\head{A finite quantitative check}
For $n=1,\ldots,8$, the exact integer coefficients and a rational enclosure for $\pi$ bound the ratio $R_n$ strictly between zero and one. In this tested range, $1-R_n<3^{-(2n+1)}$. The plot displays the logarithm of the enclosed error, not an independently asserted all-$n$ bound.
\head{Higher reciprocal powers need more Laurent data}
For positive integral order $r$, a simple denominator zero becomes a pole of order $r$. The highest Laurent term contributes a factor $\binom{m+r-1}{r-1}$, but lower Laurent terms still contribute on the same exponential scale. Subtracting \emph{all} principal parts is what leaves a holomorphic remainder controlled by the next zero radius.
\takeaway{Nearest-zero geometry explains factorial-times-exponential growth. Multiplicity, reciprocal order, and lower principal parts determine its precise normalization.}
\platecaption{The ratio certificates use a rational Machin-formula enclosure for $\pi$. The inherited general theorem retains simple nearest zeros, a gap to the next modulus, and positive integral reciprocal order. The plot does not remove those hypotheses.}

\vpage{24}{Convergence is not continuation}{ZERO ZETA FUNCTIONS}{Traditional monograph: Theorem 14.2 (p. 40), Theorem 14.3 (p. 41), Corollary 14.4 (p. 42).}
Reduced zeros grow like $m^N$. Their absolute zero-zeta series therefore has half-plane of absolute convergence $\Re s>1/N$. The absolute zero zeta in the original variable instead requires $\Re s>1$.
\vfig{zeta_plane}
For $N=2$ and $N=4$, the lattices are exact:
\[
\mathfrak Z_2^{\rm red}(s)=\pi^{-2s}\zeta(2s,1/2),\qquad
\mathfrak Z_4^{\rm red}(s)=(\sqrt2\pi)^{-4s}\zeta(4s,1/2).
\]
Each identity first holds where its defining series converges. The right side supplies its meromorphic continuation. There is a single pole, at $1/2$ or $1/4$, with residue $1/(2\pi)$ or $1/(4\sqrt2\pi)$ respectively.
\head{The factor $N$ matters twice}
For every order in the paper,
\[
\mathfrak Z_N^{\rm abs}(s)=N\mathfrak Z_N^{\rm red}(s/N).
\]
For general $N$, exponentially accurate zero locations leave an entire correction to the Hurwitz model. A finite exceptional zero set also contributes entire terms. This reasoning establishes continuation, not a zeta functional equation.
\takeaway{An analytic continuation can exist at points where the original infinite sum diverges. It must not be evaluated there by naively summing that series.}
\platecaption{The shaded regions show convergence, and the open circles mark poles of the continued functions. They are not zeros of the zeta functions. Branch-dependent orbit zetas and derivative-weighted zetas retain their separately stated conventions.}

\vpage{25}{One closed calculation: the number 69}{COEFFICIENT AND ZERO WORKED MODEL}{Traditional monograph: Corollary 7.7 (p. 16), Example 11.12 (p. 31), Theorem 12.6 (p. 36).}
Use $N=4$ throughout this page. Its denominator, its exact zero lattice, and the reciprocal recurrences provide two routes to the same coefficient. All symbols in the plate refer to this one order.
\vfig{worked_circuit}
\head{Route one: invert the denominator}
The $w^2$ coefficient of $\Delta_4^{-1}$ is
\[
\frac1{24^2}-\frac1{8!}=\frac{69}{8!}.
\]
Thus $P_{4,2}(1)=69$. Keeping $r$ variable gives $P_{4,2}(r)=35r^2+34r$.
\head{Route two: use the zero moments}
The exact reduced zeros are $4\pi^4(m+1/2)^4$. Their first two power sums are $\sigma_1=1/24$ and $\sigma_2=17/10080$. The Bell/Newton identity gives the same reciprocal coefficient $(\sigma_1^2+\sigma_2)/2$. The tangent independently reads off the second moment as $4\cdot7!\sigma_2=34$.
This links coefficients, geometry, and derivatives without a fitted parameter or an inferred zero list.
\takeaway{The content is the agreement of precisely normalized descriptions. A matching small example illustrates the general bridges. It is not a proof of all their analytic hypotheses.}
\platecaption{Every rational value and polynomial identity on this page is checked exactly. The infinite zero-sum interpretation uses the inherited canonical-product theorem and the exact order-four zero formula.}

\vpage{26}{What is evidence, and for what?}{FORMAL AND COMPUTATIONAL BOUNDARIES}{Traditional monograph: Section 15 (p. 45), Section 16 (p. 45).}
This guide accompanies the final traditional edition and the Lean guide. Its twenty-six worked pages preserve the mathematical models of the original visual appendix. Final references and evidence descriptions connect those models to the present edition.

{\small
\begin{tabularx}{\linewidth}{@{}p{.24\linewidth}Y@{}}\toprule
\textbf{Evidence layer} & \textbf{What it supports}\\\midrule
Inherited formal record & The v1.0.0 sources, theorem map, dependency pins, saved axiom reports, and original review retain their own scope. The release reports a 68-result contract. No fresh Lean compilation is claimed here.\\
Exact and symbolic suite & Fresh isolated execution passed all 147 finite, symbolic, and enclosure checks and rejected five deliberately false controls. The unchanged checker and models were run with SymPy 1.14.0 and mpmath 1.3.0.\\
Certified zero brackets & Four rational intervals of width $10^{-20}$ establish local existence and uniqueness only. The fresh full suite reproduced all four rational brackets, opposite endpoint signs, and derivative exclusion throughout each interval. These are local certificates, not a global zero census.\\
Numerical visualization & Sampled analytic curves, finite products, moment shares, and displacement diagnostics have the roles stated on their pages. They are not exact global zero censuses.\\
Figure provenance & Fresh isolated model-driven regeneration reproduced all twenty-four editable TikZ files. Twenty-three remain identical to the original release. On V13, one label was moved down to clear the loop arrow. The mathematical labels and objects are unchanged.\\
Editorial preservation & The original worked formulas, model code, checker, and vector figure sources are preserved. Updated locators refer to the final traditional edition.\\\bottomrule
\end{tabularx}}

\head{Boundaries that remain in force}
The general family has eventual positivity and simplicity of its reduced zeros, not an assertion that every reduced zero is positive and simple. The paper leaves universal simplicity, positivity of every reduced zero, full displacement expansions, and a zeta functional equation open. The exact order-two and order-four models do not settle those questions for every order.
Imported standard fact 11.1 stays outside the 68-result contract. Fixed logarithms, zero-free domains, multiplicities, and the hypotheses of nearest-pole asymptotics remain explicit. No historical priority is inferred from successful computation or formal verification.

\head{Replay the worked layer}
Run \texttt{python build.py} in this editable guide folder to compile the PDF in an isolated temporary directory. The separate model checker is \texttt{checks/verify\_examples.py}. The plate generator is \texttt{source/generate\_graphics.py}. They require the listed Python dependencies. Run those mathematical tools in a temporary copy to preserve the delivered evidence.

\takeaway{A proof, a finite exact test, a root enclosure, and a sampled curve answer different questions. Keeping their scopes explicit makes the pictures useful mathematical evidence.}
\platecaption{The source package preserves the formal release, proofread review, independent audit, exact certificates, and rebuild instructions. PDF assembly is a separate operation from mathematical verification.}

\endgroup

\end{document}
